\magnification = 1200
\input epsf
\baselineskip = 20pt
\voffset = .75truein
\hsize = 6.5truein
\vsize = 8.5truein
\def\bul{$\bullet$}
\def\cir{$\circ$}
\def\no{\noindent}
\def\newpage{\vfill\eject}
\font\smcaps=cmcsc10
\overfullrule=0pt
\ContentsSwitchtrue

\footline={\hfil}

{

\parskip=0pt
\baselineskip=13pt

\no Federal Reserve Bank of Minneapolis\par
\no Research Department Staff Report 362
\medskip
\smallskip
\no Revised December 2006\par
\vskip 1.0truein

\no{\twelvepoint\bf APPENDICES:
Business Cycle Accounting${}^\dagger$}

\bigskip
\smallskip
\no{\smcaps V.V.~Chari}\nobreak\par
\no{University of Minnesota}\nobreak\par
\no{and Federal Reserve Bank of Minneapolis}

\bigskip
\smallskip
\no{\smcaps Patrick J.~Kehoe}\nobreak\par
\no{Federal Reserve Bank of Minneapolis}\nobreak\par
\no{and University of Minnesota}

\bigskip
\smallskip
\no{\smcaps Ellen R.~McGrattan}\nobreak\par
\no{Federal Reserve Bank of Minneapolis}\nobreak\par
\no{and University of Minnesota}

\vskip 3.5truein

{
\baselineskip=13pt

\no ${}^\dagger$
The views expressed herein are those of the authors and not necessarily
those of the Federal Reserve Bank of Minneapolis or the Federal Reserve
System.

}
}

\newpage
\footline={\hss\tenrm\folio\hss}
\pageno=-1
{

\baselineskip = 15pt

\nosechead{Table of Contents}
\tag{Pg.tabcontents}{\folio}
\TOCmargin=0cm
\contents

}


\newpage
\footline={\hss\tenrm\folio\hss}
\pageno=1

\Appendix{A}{Computation and Estimation}

\noindent
In this Appendix, we describe the prototype model 
used in ``Business Cycle Accounting''  
and the details of the computation of equilibria
and estimation of parameters.
We also discuss the sensitivity analysis 
we did for the benchmark
prototype model allowing for variable capital utilization,
investment adjustment costs, and an alternative 
specification for the investment wedge.
Finally, we include figures, tables, and proofs
not shown in the main text. 

\bigskip
\section{The Benchmark Model}

\noindent Below we will use the following notation for
our model variables:
\item{} $N$: population ($N_t=(1+g_n)^t$)
\item{} $c$: per-capita consumption 
\item{} $x$: per-capita investment 
\item{} $k$: per-capita net capital stock
\item{} $l$: per-capita labor input
\item{} $t\!r$: per-capita government transfers
\item{} $C$: total consumption  ($C_{t}=N_t c_t$) 
\item{} $X$: total investment 
\item{} $K$: total stock of capital
\item{} $L$:  total labor input in production
\item{} $Z$:  labor-augmenting technical change ($Z_t=z_t(1+g_z)^t$)
\item{} $r$: rental rate on capital
\item{} $w$: wage rate 
\item{} $\tau_v$: tax rate on $v$ 
\item{} $\hat v$: detrended per-capita variable $V$ 
(i.e., $\hat v_t=V_t/[N_t(1+g_z)^t]$)

\bigskip
Consider an economy with households,
firms, and the government.
The representative household chooses consumption,
investment, and labor to solve the following maximization
problem:
$$\EQNalign{
 & \max_{\{c_t,x_t,l_t\}} E\,\sum_{t=0}^\infty \beta^t 
          \, U(c_t,1-l_t)  N_t                    \EQN utility\cr
{\rm subj.\ to}\ \ 
 &  (1+\tau_{ct})c_t+ (1+\tau_{xt}) x_t 
       = (1-\tau_{kt})r_t k_t + (1-\tau_{lt}) w_t l_t
                   +\tau_{kt}\delta k_t + t\!r_t  \EQN budget\cr
 & N_{t+1} k_{t+1} = [(1-\delta)k_t + x_t] N_t    \EQN capacc\cr
 & c_t,x_{t} \geq 0 \quad {\rm in\  all\  states}, \EQN nonneg\cr
}$$
taking processes for the rental rate, wage rate, the tax rates,
and transfers as given.
The representative firm solves a simple static problem at $t$:
$$
\max_{\{K_t,L_t\}}\  
   F(K_t,Z_t L_t) - r_t K_t - w_t L_t.
$$
The government sets rates of taxes and transfers in such
a way that their budget constraint at $t$, namely,
$$
  G_t + N_tt\!r_t = \tau_{kt}(r_t-\delta) N_tk_t + \tau_{lt}w_tl_t N_t
             +\tau_{ct}N_tc_t+\tau_{xt}N_tx_t,
$$
is satisfied.
In equilibrium, the following conditions must hold:
$$\eqalign{
& N_t(c_t+x_t) +G_t 
               = F(K_t,Z_tL_t)\cr
& N_t k_t = K_t \cr
& N_t l_t = L_t. \cr
}
$$

We now derive first-order conditions in this economy.
The Lagrangian for the household optimization problem is given
by
$$\eqalign{
{\cal L} &= E\sum_t \tilde \beta^t \, 
     \biggl\{
         U(\hat c_t,1-l_t)+ {\zeta\over 3} \min(\hat x_t,0)^3 \cr
     \noalign{\medskip}
   & \quad
 + \mu_t \Big\{(1-\tau_{kt})r_t \hat k_t + (1-\tau_{lt}) 
    \hat w_t l_t + \tau_{kt}\delta \hat k_t + \hat t\!r_t
    -  (1+\tau_{ct})\hat c_t-(1+\tau_{xt})\hat x_t\Big\}\cr
     \noalign{\medskip}
   & \quad
 + \lambda_t\Big\{ (1-\delta)\hat k_t + \hat x_t 
                       - (1+g_z)(1+g_n)\hat k_{t+1}\Big\}\biggr\},\cr
}
$$
where $\tilde \beta=\beta (1+g_n) h(1+g_z)$ and $h(\cdot)$
depends on our choice of utility. If $U(c,1-l)=c^{1-\sigma}v(l)$,
then $h(1+g_z)=(1+g_z)^{1-\sigma}$.
Notice that the Lagrangian has no term for the nonnegativity
constraint on investment.
Instead, we have included a penalty function indexed
by $\zeta$.
As $\zeta$ approaches infinity, the solution to the problem
with a penalty function and no constraint on investment
is the same as the solution to the original problem with $\zeta=0$
and $x_t\geq 0$ imposed.

The first-order conditions for the problem are
$$\EQNalign{  
 & {U_2(\hat c_t,1-l_t)\over U_1(\hat c_t,1-l_t)}
   =  { 1-\tau_{lt}\over 1+\tau_{ct}} \hat w_t \cr
     \noalign{\medskip}
 & {1+\tau_{xt}\over 1+\tau_{ct}}
    U_1(\hat c_t,1-l_t)-\zeta\min(\hat x_t,0)^2 \cr
     \noalign{\medskip}
 &\qquad\qquad\qquad
   = \hat \beta
     E_t\biggl[{U_1(\hat c_{t+1},1-l_{t+1})\over
                  1+\tau_{ct+1} }\cr
     \noalign{\medskip}
 &\qquad\qquad\qquad\qquad\qquad
             \bigl\{ (1-\tau_{kt+1})r_{t+1}+\delta \tau_{kt+1}
                                    +(1-\delta)(1+\tau_{xt+1})\bigr\}\cr
     \noalign{\medskip}
 &\qquad\qquad\qquad\qquad\qquad\qquad\qquad
              -(1-\delta)\zeta \min(\hat x_{t+1},0)^2
               \biggr], \EQN dynamic \cr
}
$$
where $\hat \beta=\beta h(1+g_z)/(1+g_z)$. If $U(c,l)=c^{1-\sigma}v(l)$,
then $\hat \beta = \beta(1+g_z)^{-\sigma}$.

In addition, we have first-order conditions for the firm's 
static problem.  These are
$$\eqalign{
r_t  &= F_1(\hat k_t,z_t l_t)\cr
\hat w_t  &= F_2(\hat k_t,z_t l_t)z_t.\cr
}
$$
Finally,  we have a resource constraint given by
$$\hat c_t + \hat g_t +\hat x_t
               = F(\hat k_t,z_tl_t)$$
once we detrend variables.


\bigskip
\section{Algorithms for Computing Equilibria}

Below we show how to compute the equilibrium
using a nonlinear method and a log-linear method.
We need to use the nonlinear methods to compute
equilibrium paths during the 1930s because the 
declines in aggregate variables are so large.
We use the log-linear method when we 
derive estimates of the process for the wedges
and for computing equilibria in the postwar period.
The log-linear method is convenient with our maximum likelihood
estimation (MLE) procedure because a nonlinear method
with a very large state space is computationally demanding
when computing likelihood estimates.\note{%
We experimented with very different expectations and 
found that they had only a very small impact on our results 
and, therefore, did not affect our conclusions.}  
With estimates of our stochastic process
we determine expectations.  These expectations are
inputs to our nonlinear model and used when we 
do our accounting exercise.

From here on, we make the following functional form
assumptions and auxiliary choices:
\item{} $F(k,l) = k^\theta l^{1-\theta}$
\item{} $U(c,1-l) = (c (1-l)^\psi)^{1-\sigma}/(1-\sigma)$
\item{} $s_{t+1} = P_0 + P s_t + Q \epsilon_{s,t+1},\ \  
               \epsilon_s\sim N(0,I)$
         (which is represented either as a Markov chain using the method
         of Tauchen 1986 or as a continuous process).
\item{} $\log(z_t)= \log z(s_t)$
\item{} $\log \hat g_t= \log \hat g(s_t)$
\item{} $\tau_{lt}= \tau_l(s_t)$
\item{} $\tau_{xt}= \tau_x(s_t)$
\item{} $\tau_{kt}= \tau_k(s_t)$
\item{} $\tau_{ct}= \tau_c(s_t)$

\bigskip
\bigskip
\subsection{Nonlinear Computation}

For our nonlinear solution method, we assume that the
vector autoregressive process for the state can be well approximated
by a  Markov chain.  
Let $s$ be the index for the state.  Then at time $t$, if the state
is $s$,
$\hat g_t= \hat g(s)$,
$\tau_{lt} = \tau_l(s)$,
$\tau_{xt} = \tau_x(s)$, and
$z_t =  z(s)$.
The transition matrix for $s$ is given by $\Pi$
with $\pi(s,s')$ being the probability of going from
state $s$ to state $s'$.

The state of the economy in any period
can be summarized by two scalars: $\hat k$ and $s$.
Our Fortran code computes
the decision rule $\hat c(\hat k,s)$.  All other
decisions can be determined via static first-order conditions
once we know $\hat c(\hat k,s)$.
In particular,
$l(\hat k,s)$ and $\hat x(\hat k,s)$ can 
be determined once we know consumption.

To compute $\hat c(\hat k,s)$, we apply the finite-element method
using the dynamic first-order condition as the residual
and Galerkin bases.
More specifically, we
assume that the consumption function is well approximated by
$$ \hat c(\hat k,s) = \sum_{j=1}^{nnodes} \alpha_j^s \Psi_j(k),$$
where the $\Psi_j$ is a function that takes on
nonzero values in two cells (or ``elements'')
of a grid over $\hat k$ around grid
point (or ``node'') $j$.
The algorithm is to find the coefficients $\alpha_j^s $,
$j=1,\ldots nnodes$, $s=1,\ldots S$
that satisfy the following equations:
$$ \int R(\hat k,s;\alpha) \Psi_j(\hat k) \, d\hat k =0
$$
for all $s$ and $j$ where
$$\eqalign{
R( & \hat k,s;\alpha) =
  {1+\tau_x(s)\over 1+\tau_c(s)} U_1(\hat c,1-l)\cr
   &\qquad\qquad
    -\zeta \min(\hat x,0)^2
  +\hat \beta(1-\delta)
        \zeta \sum_{s'}\pi_{s,s'} \min (\hat x',0)^2 \cr
\noalign{\medskip}
& \qquad\qquad
  - \hat \beta \sum_{s'} \pi_{s,s'} \,
  {U_1(\hat c',1-l')\over 1+\tau_c(s')}\cr
& \qquad\qquad\qquad
 \biggl\{(1-\tau_k(s'))F_1(\hat k',z(s')l')+\tau_k(s')\delta
    +(1-\delta)(1+\tau_x(s))
          \biggr\}. \cr
}
$$
The investments $x$ and $x'$ satisfy resource constraints
$$\eqalign{
  \hat x &= F(\hat k,z(s)l)-\sum_j\alpha_j^s \Psi_j(\hat k)
             -\hat g(s)\cr
        \noalign{\medskip}
  \hat x'&= F(\hat k',z(s')l')-\sum_j\alpha_j^s \Psi_j(\hat k')
             -\hat g(s').\cr
}
$$
The next period capital stock is given by
$$\eqalign{
\hat k' & = ((1-\delta) \hat k + \hat x)/[(1+g_n)(1+g_z)]. \cr
}
$$
The labor inputs $l$ and $l'$ solve
$$\eqalign{
&
{U_2(\hat c,1-l)
  \over U_1(\hat c,1-l)} =
 {1-\tau_l(s) \over
         1+\tau_c(s)}
         F_2(\hat k,z(s)l)z(s) \cr
\noalign{\bigskip}
&
{U_2(\hat c',1-l') \over
   U_1(\hat c',1-l') }=
{1-\tau_l(s')\over 1+\tau_c(s')} F_2(\hat k',z(s')l')z(s').\cr
}
$$


\bigskip
\bigskip
\subsection{Log-Linear Computation}

We now describe the steps taken for the log-linear solution method
(with an interior solution and $\zeta=0$).
Because we are going to apply maximum likelihood estimation,
we will derive the solution analytically.

We start by writing the system of equations in terms of $k$ and $s$.
This is done by replacing $r$, $w$, $\hat c$, and $\hat x$
in the first-order conditions with functions of the states. Thus
we start with
$$\EQNalign{
 &\hat c_t + \hat g_t +(1+g_z)(1+g_n)\hat k_{t+1}-(1-\delta)\hat k_t
    = \hat y_t = \hat k_t^\theta(z_tl_t)^{1-\theta}
\EQN {foc1}\cr
 &{\psi \hat c_t\over 1-l_t}
    = (1-\tau_{lt})(1-\theta) \hat k_t^\theta l_t^{-\theta} z_t^{1-\theta}
\EQN {foc2}\cr
 &(1+\tau_{xt}) \hat c_t^{-\sigma} (1-l_t)^{\psi (1-\sigma)}
     \cr
 &\qquad
       = \hat \beta E_t
            \hat c_{t+1}^{-\sigma} (1-l_{t+1})^{\psi (1-\sigma)}
          [\theta \hat k_{t+1}^{\theta-1} (z_{t+1}l_{t+1})^{1-\theta}
              +(1-\delta)(1+\tau_{xt+1})],
\EQN {foc3}\cr
}
$$
which can be reduced to the following:
$$\eqalign{
 &\psi [\hat k_t^\theta(z_tl_t)^{1-\theta}
             -(1+g_n)(1+g_z)\hat k_{t+1}+(1-\delta)\hat k_t-\hat g_t]\cr
 &\qquad\qquad
    = (1-\tau_{lt})(1-\theta) \hat k_t^\theta l_t^{-\theta} 
          z_t^{1-\theta}(1-l_t) \cr
 &(1+\tau_{xt}) [\hat k_t^\theta(z_tl_t)^{1-\theta}
             -(1+g_n)(1+g_z)\hat k_{t+1}+(1-\delta)\hat k_t-\hat g_t]^{-\sigma} 
              (1-l_t)^{\psi (1-\sigma)}
     \cr
 &\qquad\qquad
       = \hat \beta E_t[\hat k_{t+1}^\theta(z_{t+1}l_{t+1})^{1-\theta}
             -(1+g_n)(1+g_z)\hat k_{t+2}\cr
 &\qquad\qquad\qquad\qquad
    +(1-\delta)\hat k_{t+1}-\hat g_{t+1}]^{-\sigma} 
                (1-l_{t+1})^{\psi (1-\sigma)}\cr
 &\qquad\qquad\qquad\qquad
          [\theta \hat k_{t+1}^{\theta-1} (z_{t+1}l_{t+1})^{1-\theta}
              +(1-\delta)(1+\tau_{xt+1})].\cr
}
$$

Next, we compute the steady state of the system for constant
values for $z$, the taxes, and government spending:
$$\eqalign{
  & \hat k/l = \left({(1+\tau_x)(1-\hat\beta(1-\delta)) \over
                  \hat \beta \theta z^{1-\theta}}\right)^{1/(\theta-1)}\cr
  & \hat c = \left[(\hat k/l)^{\theta-1}z^{1-\theta}
              -(1+g_z)(1+g_n)+1-\delta\right] \hat k-\hat g =  
                      \xi_1 \hat k -\hat g\cr
  &\hat c = 
     \left[(1-\tau_l)(1-\theta) (\hat k/l)^{\theta}z^{1-\theta}/\psi\right]
             (1- 1/(\hat k/l) \ \hat k)
         = \xi_2-\xi_3\hat k,\cr
}
$$
where the last two equations imply
$\hat k=(\xi_2+\hat g)/(\xi_1+\xi_3)$, $\hat c=\xi_1 \hat k -\hat g$, 
$l=(1/(\hat k/l)) \hat k$.

The log-linearization is done around these steady-state values.
Detrended consumption is given approximately by
$$\eqalign{
  \hat c_t &\approx \hat c \log \hat c_t\cr
           &\approx \hat k^{\theta} (zl)^{1-\theta}
           [\theta \log \hat k_t + (1-\theta)(\log z_t+\log l_t)]\cr
          &\quad
          -(1+g_z)(1+g_n)\hat k \log \hat k_{t+1} 
          +(1-\delta)\hat k \log \hat k_t-\hat g \log\hat g_t. \cr
}
$$
The labor input is then derived from the static first-order condition
\Ep{foc2}:
$$\eqalign{
    0 &\approx \psi\bigl\{\hat k^{\theta} (zl)^{1-\theta}
           [\theta \log \hat k_t + (1-\theta)(\log z_t+\log l_t)]\cr
          &\qquad\quad
          -(1+g_z)(1+g_n)\hat k \log \hat k_{t+1} 
          +(1-\delta)\hat k \log \hat k_t-\hat g\log \hat g_t\bigr\} \cr
          &\qquad\qquad\quad 
                 + (1-\theta)(1-\tau_l)\hat k^\theta 
                          l^{-\theta} z^{1-\theta}(1-l)
                    \bigl\{1/(1-\tau_l)\,\tau_{lt}\cr
          &\qquad\qquad\qquad\quad 
                               -\theta \log \hat k_t
                                  +\theta\log l_t -(1-\theta) \log z_t
                                  +l/(1-l) \log l_t \bigr\},\cr
}
$$
which we write succinctly as
$$
      \log l_t = \phi_{lk}\log \hat k_t
            +\phi_{lz}\log z_t
            +\phi_{ll}\tau_{lt}
            +\phi_{lg}\log \hat g_t
            +\phi_{lk'}\log \hat k_{t+1}.
\EQN logl
$$
Using this equation for $\log l$, we use the other
static first-order conditions to write 
$\log \hat y$,
$\log \hat x$, and
$\log \hat c$  as follows:
$$\EQNalign{
 \log \hat y_t & = \phi_{yk}\log \hat k_t
            +\phi_{yz}\log z_t
            +\phi_{yl}\tau_{lt}
            +\phi_{yg}\log \hat g_t
            +\phi_{yk'}\log \hat k_{t+1} \cr
      &= (\theta+(1-\theta)\phi_{lk})\log \hat k_t
                        +(1-\theta)(1+\phi_{lz}) \log z_t \cr
      &\qquad\qquad\qquad\qquad
                        +(1-\theta)[\phi_{ll}\tau_{lt}
                                 +\phi_{lk'}\log \hat k_{t+1}]
\EQN loghaty\cr
 \log \hat x_t &= (1+g_z)(1+g_n) \hat k/\hat x \log \hat k_{t+1}
                    -(1-\delta) \hat k/\hat x \log \hat k_t
\EQN loghatx\cr
 \log \hat c_t &= \phi_{ck}\log \hat k_t
            +\phi_{cz}\log z_t
            +\phi_{cl}\tau_{lt}
            +\phi_{cg}\log \hat g_t
            +\phi_{ck'}\log \hat k_{t+1} \cr
      & = [\hat y\log y_t
                      -\hat x\log x_t -\hat g\log \hat g_t]/\hat c,
\EQN loghatc\cr
}
$$
where the $\phi$'s are known functions of the parameters.

Capital is derived from the dynamic first-order condition
$$ \EQNalign{
0&\approx
  (1+\tau_x)\hat c^{-\sigma} (1-l)^{\psi(1-\sigma)} 
        \left\{-\psi(1-\sigma) l/(1-l) \log l_t
                  -\sigma \log \hat c_t \right\} \cr
 & \qquad \qquad\qquad
   +\hat c^{-\sigma} (1-l)^{\psi(1-\sigma)} \tau_{xt} \cr
 & \qquad 
   -\hat\beta 
        E_t\bigl\{ \bigl[\theta \hat k^{\theta-1}(zl)^{1-\theta}
           +(1-\delta)(1+\tau_x)\bigr]\cr
 & \qquad \qquad\qquad
       \cdot\bigl[\hat c^{-\sigma} (1-l)^{\psi(1-\sigma)}
            \left\{-\psi(1-\sigma) l/(1-l) \log l_{t+1}
                    -\sigma \log \hat c_{t+1} \right\}
       \bigr]\cr
 & \qquad \qquad\qquad
        +\hat c^{-\sigma} (1-l)^{\psi(1-\sigma)}
             \bigl[\theta \hat k^{\theta-1}(zl)^{1-\theta}(1-\theta)
        \cr
 & \qquad \qquad\qquad\qquad\qquad
                  \cdot(\log l_{t+1} + \log z_{t+1}-\log \hat k_{t+1})
                  +(1-\delta)\tau_{xt+1}\bigl]\bigr\},
\EQN dynamic\cr
}
$$
which simplifies to
$$ \eqalign{
0&\approx
  (1+\tau_x)
        \left\{ -\psi(1-\sigma) l/(1-l) \log l_t
                -\sigma \log \hat c_t \right\} 
   +\tau_{xt} \cr
 & \qquad 
       -E_t\bigl\{ (1+\tau_x)
            \left\{-\psi(1-\sigma) l/(1-l) \log l_{t+1}
                   -\sigma \log \hat c_{t+1}
            \right\}\cr
 & \qquad \quad
     +\hat\beta\bigl[r(1-\theta)
                  (\log l_{t+1} + \log z_{t+1}-\log \hat k_{t+1})
                  +(1-\delta)\tau_{xt+1}\bigr]\bigr\},\cr
}
$$
where $r=\theta \hat y/\hat k$.
          
We guess the following form of the solution for capital:
$$ \log(\hat k_{t+1}) = \gamma_0 + \gamma_k \log \hat k_t
                         + \gamma_z \log z_t
                 + \gamma_l \tau_{lt}
                 +\gamma_x \tau_{xt}
                 +\gamma_g \log\hat g_t
\EQN{formofsolution}
$$
and set $\gamma$'s so that the dynamic residual \Ep{dynamic} is 
exactly 0. We 
can do this by first ignoring shock terms and find $\gamma_k$ that
satisfies a particular quadratic equation (that does not depend
on any other $\gamma$ coefficient).  Then,
we can find $\gamma_z$, $\gamma_l$, $\gamma_x$, and $\gamma_g$ by solving
a linear system of equations with $\gamma_k$ assumed known.
Finally, we use the steady-state equations to determine $\gamma_0$.

We start by deriving $\gamma_k$.  To do this, 
we need to write out the coefficients on $\hat k_{t+2}$, 
$\hat k_{t+1}$, and $\hat k_t$ in the dynamic first-order condition
\Ep{dynamic} making use
of the $\phi$'s from the static first-order conditions 
\Ep{logl}--\Ep{loghatc}.
For now, we 
can ignore the expectations operator.  We get the following:
$$ \eqalign{
0&\approx
  -(1+\tau_x)
        \bigl\{ -\psi(1-\sigma) l/(1-l)
              [\phi_{lk}\log \hat k_t+\phi_{lk'} \log \hat k_{t+1}] 
   \cr
 & \qquad\qquad\qquad \qquad
               -\sigma 
              [\phi_{ck}\log \hat k_t+\phi_{ck'} \log \hat k_{t+1}] 
         \bigr\} \cr
 & \quad 
       +E_t(1+\tau_x)
            \bigl\{-\psi(1-\sigma) l/(1-l) 
              [\phi_{lk}\log \hat k_{t+1}+\phi_{lk'} \log \hat k_{t+2}] \cr
 & \qquad\qquad\qquad\qquad\qquad \qquad
               -\sigma 
              [\phi_{ck}\log \hat k_{t+1}+\phi_{ck'} \log \hat k_{t+2}] 
            \bigr\}\cr
 & \qquad \qquad\quad
     +\hat\beta E_t \bigl[r(1-\theta)
                  ((\phi_{lk}-1)\log \hat k_{t+1}+\phi_{lk'}\hat k_{t+2})
              \bigr],\cr
}
$$
which simplifies to 
$$\EQNalign{
  0 = &[\hat\beta r(1-\theta)\phi_{lk'}
                 -(1+\tau_x)\psi(1-\sigma)l/(1-l)\phi_{lk'}
         -(1+\tau_x)\sigma\phi_{ck'}] \hat k_{t+2}\cr
      +&[\hat\beta r(1-\theta)(\phi_{lk}-1)
               -(1+\tau_x)\psi(1-\sigma)l/(1-l)\phi_{lk}
               -(1+\tau_x)\sigma\phi_{ck}
       \cr
       &\qquad\qquad\qquad\qquad
               +(1+\tau_x)\psi(1-\sigma)l/(1-l)\phi_{lk'}
               +(1+\tau_x)\sigma\phi_{ck'}
       ] \hat k_{t+1}\cr
      -&[-(1+\tau_x)\psi(1-\sigma)l/(1-l)\phi_{lk}
               -(1+\tau_x)\sigma\phi_{ck}
       ] \hat k_t\cr
      +&{\rm \ all\ other\ terms}
\EQN dynamicpiece
}
$$
or, more succinctly, rewrite \Ep{dynamicpiece} as
$ (a+ bL + cL^2) \hat k_{t+2} =$ other terms,
where $\gamma_k$ is the root of the quadratic 
inside the unit circle.  Note that $\gamma_k$ does not depend
on the other unknown $\gamma$'s.

Given $\gamma_k$, we can solve a linear system
for the other $\gamma$'s.  At this point,
we do not ignore the expectations operator:
$$ \eqalign{
0&\approx
  (1+\tau_x)
        \bigl\{-\psi(1-\sigma) l/(1-l)
              [(\phi_{lz}+\phi_{lk'}\gamma_z)\log z_t\cr
 & \qquad\qquad\qquad \qquad\qquad\qquad\qquad
               +(\phi_{ll}+\phi_{lk'}\gamma_l)\tau_{lt}\cr
 & \qquad\qquad\qquad \qquad\qquad\qquad\qquad
               +\phi_{lk'}\gamma_x\tau_{xt} 
               +(\phi_{lg}+\phi_{lk'}\gamma_g)\log\hat g_t\cr
 & \qquad\qquad\qquad \qquad
               -\sigma 
              [(\phi_{cz}+\phi_{ck'}\gamma_z)\log z_t
               +(\phi_{cl}+\phi_{ck'}\gamma_l) \tau_{lt}\cr
 & \qquad\qquad\qquad \qquad\qquad\qquad
               +\phi_{ck'}\gamma_x \tau_{xt}
               +(\phi_{cg}+\phi_{ck'}\gamma_l) \log\hat g_t
               ]\bigr\} +\tau_{xt}\cr
}
$$
$$\eqalign{
 & \qquad 
       -E_t(1+\tau_x)
            \bigl\{-\psi(1-\sigma) l/(1-l) 
              [(\phi_{lk}+\phi_{lk'}\gamma_k)(\gamma_z\log z_t
                                              +\gamma_l\tau_{lt}
                                              +\gamma_x\tau_{xt}
                                              +\gamma_g\log\hat g_t)\cr
 & \qquad\qquad\qquad \qquad\qquad\qquad\qquad \qquad
              + (\phi_{lz}+\phi_{lk'}\gamma_z)\log z_{t+1}\cr
 & \qquad\qquad\qquad \qquad\qquad\qquad\qquad \qquad
              + (\phi_{ll}+\phi_{lk'}\gamma_l)\tau_{lt+1}\cr
 & \qquad\qquad\qquad \qquad\qquad\qquad\qquad \qquad
              + \phi_{lk'}\gamma_x\tau_{xt+1}\cr
 & \qquad\qquad\qquad \qquad\qquad\qquad\qquad \qquad
              + (\phi_{lg}+\phi_{lk'}\gamma_g)\log\hat g_{t+1}]\cr
}
$$
$$\eqalign{
 & \qquad\qquad\qquad\qquad\qquad \qquad
               -\sigma 
              [(\phi_{ck}+\phi_{ck'}\gamma_k)(\gamma_z\log z_t
                                              +\gamma_l\tau_{lt}
                                              +\gamma_x\tau_{xt}
                                              +\gamma_g\log\hat g_t)\cr
 & \qquad\qquad\qquad\qquad\qquad \qquad\qquad
               +(\phi_{cz}+\phi_{ck'}\gamma_z)\log z_{t+1}\cr
 & \qquad\qquad\qquad\qquad\qquad \qquad\qquad
               +(\phi_{cl}+\phi_{ck'}\gamma_l)\tau_{lt+1}\cr
 & \qquad\qquad\qquad\qquad\qquad \qquad\qquad
               +\phi_{ck'}\gamma_x\tau_{xt+1}\cr
 & \qquad\qquad\qquad\qquad\qquad \qquad\qquad
               +(\phi_{cg}+\phi_{ck'}\gamma_g)\log\hat g_{t+1}]\bigr\}\cr
}
$$
$$\eqalign{
 & \qquad \qquad\quad
     -\hat\beta r(1-\theta) E_t \bigl[(\phi_{lk}+\phi_{lk'}\gamma_k-1)
                    (\gamma_z\log z_t + \gamma_l\tau_{lt}+\gamma_x\tau_{xt}
                      +\gamma_g\log\hat g_t)\cr
 & \qquad\qquad\qquad\qquad\qquad \qquad\qquad
                    +(\phi_{lz}+\phi_{lk'}\gamma_z+1)\log z_{t+1}\cr
 & \qquad\qquad\qquad\qquad\qquad \qquad\qquad
                    +(\phi_{ll}+\phi_{lk'}\gamma_l)\tau_{lt+1}\cr
 & \qquad\qquad\qquad\qquad\qquad \qquad\qquad
                    +\phi_{lk'}\gamma_x\tau_{xt+1}\cr
 & \qquad\qquad\qquad\qquad\qquad \qquad\qquad
                    +(\phi_{lg}+\phi_{lk'}\gamma_g)\log\hat g_{t+1}\bigr]\cr
 & \qquad\qquad\qquad
     -\hat \beta (1-\delta) E_t \tau_{xt+1},\cr
}
$$
which reduces to
$$\EQNalign{
 0 &= (\kappa_0+\kappa_1\cdot\gamma_s)' 
           [\log z_t,\tau_{lt},\tau_{xt},\log\hat g_t]'
     + (\zeta_0+\zeta_1\cdot\gamma_s)' 
          E_t [\log z_{t+1},\tau_{lt+1},\tau_{xt+1},\log\hat g_{t+1}]'
    \cr
   &= (\kappa_0+\kappa_1\cdot\gamma_s)' 
           [\log z_t,\tau_{lt},\tau_{xt},\log\hat g_t]'
     + (\zeta_0+\zeta_1\cdot\gamma_s)' P
           [\log z_t,\tau_{lt},\tau_{xt},\log\hat g_t]',\EQN lqsystem
    \cr
}
$$
where $x\cdot y$ denotes element-by-element multiplication
of vectors $x$ and $y$;
$\kappa_0$, $\kappa_1$, $\zeta_0$, and $\zeta_1$ are vectors of 
length 4 with elements equal to functions
of the parameters and $\gamma_k$; and 
$\gamma_s=[\gamma_z,\gamma_l,\gamma_x,\gamma_g]'$.
The $\gamma$'s that set this residual to zero satisfy a four-dimensional
linear system,
$$
 (\kappa_0+\kappa_1\cdot\gamma_s)' 
     + (\zeta_0+\zeta_1\cdot \gamma_s) ' a = 0.
$$


\bigskip
\bigskip
\subsubsection{Checking a Test Case}

Assume $g_n=g_z=\psi=0$, $\sigma=\delta=1$, 
$\hat g_t=\tau_{lt}=\tau_{xt}=0$, and
$$\log z_{t+1}=\rho_0+ \rho \log z_t + \epsilon_{z,t+1}
$$
so that
$$ P   =\left[\matrix{\rho_z & 0 & 0 & 0 & \rho_0\cr
                           0 & 0 & 0 & 0 & 0\cr
                           0 & 0 & 0 & 0 & 0\cr
                           0 & 0 & 0 & 0 & 0\cr
                           0 & 0 & 0 & 0 & 1\cr}\right],\quad
   Q   =\left[\matrix{\sigma_z & 0 & 0 & 0 \cr
                        0      & 0 & 0 & 0 \cr
                        0      & 0 & 0 & 0 \cr
                        0      & 0 & 0 & 0 \cr
                        0      & 0 & 0 & 0 \cr}\right].
$$
The solution in this case: 
$$ \log \hat k_{t+1} = \log(\beta\theta) + \theta \log \hat k_t
                     +(1-\theta)\log z_t.
$$

Using the formulas above, we have the following steady state:
$$ \eqalign{
   \hat k &= (\beta\theta)^{1/(1-\theta)} z\cr
   \hat c &= (1-\beta\theta)/(\beta\theta)\hat k\cr
   \hat y &= z^{1-\theta}\hat k^\theta   \cr
}
$$
and $\gamma_k$ from quadratic
$$\eqalign{
  0 = &-\phi_{ck'} \hat k_{t+2}
      -[\hat\beta r(1-\theta)+\phi_{ck}-\phi_{ck'}] \hat k_{t+1}
      +\phi_{ck}\hat k_t+{\rm all\ other\ terms}\cr
    = &\beta\theta/(1-\beta\theta) \hat k_{t+2}
      -[(1-\theta)+\theta/(1-\beta\theta) 
           -\beta\theta/(1-\beta\theta)] \hat k_{t+1}
      +[\theta/(1-\beta\theta)] \hat k_t\cr
    = &\beta\theta\hat k_{t+2}
      -[(1-\theta)(1-\beta\theta)+\theta-\beta\theta] \hat k_{t+1}
      +\theta \hat k_t,\cr
}
$$
which implies $\gamma_k=\theta$.
The other coefficients satisfy
$$ \eqalign{
0&\approx
        -(\phi_{cz}+\phi_{ck'}\gamma_z)\log z_t
               -\phi_{ck'}\gamma_l\tau_{lt}
               -\phi_{ck'}\gamma_x\tau_{xt} 
               -\phi_{ck'}\gamma_g\log\hat g_t +\tau_{xt}\cr
 & \qquad 
       +(\phi_{ck}+\phi_{ck'}\theta)(\gamma_z\log z_t
                                              +\gamma_l\tau_{lt}
                                              +\gamma_x\tau_{xt}
                                              +\gamma_g\log\hat g_t)\cr
 & \qquad\qquad
           +(\phi_{cz}+\phi_{ck'}\gamma_z) E_t\log z_{t+1}\cr
 & \qquad \qquad\quad
     +(1-\theta) \bigl[\gamma_z\log z_t+\gamma_l\tau_{lt}
                                 +\gamma_x\tau_{xt}
                                 +\gamma_g\log\hat g_t
                            -E_t\log z_{t+1}\bigr]\cr
}
$$
and, therefore, 
$\gamma_z=1-\theta$, $\gamma_l=\gamma_g=0$, $\gamma_x=-1+\beta\theta$.


\bigskip
\bigskip
\subsubsection{Computation for the General Log-Linear Case}

Assume that the solution is
$$ \log \hat k_{t+1}
   = a \log \hat k_t + b
        \left[\matrix{ \log z_t & \tau_{lt} & \tau_{xt} & \log\hat g_t}
        \right]'+{\rm constant,}
$$
where $a$ is a scalar and $b$ is $1\times 4$.
Assume that the residual from the dynamic first-order condition is
$$\eqalign{
& f(E_t \log \hat k_{t+2}, \log\hat k_{t+1}, \log \hat k_t,
             \log z_{t+1}, \log z_t, \tau_{lt+1}, \tau_{lt}, 
             \tau_{xt+1}, \tau_{xt}, \log\hat g_{t+1}, \log\hat g_t)\cr
&\qquad   \approx
  a_0 E_t \log \hat k_{t+2}+ a_1 \log\hat k_{t+1} +a_2 \log \hat k_t
          + b_0 E_t s_{t+1} + b_1 s_t.\cr
}
$$     
Then the general solution algorithm is to find 
$a$ that solves the quadratic equation
$$  a_0 a^2+a_1 a+a_2= 0$$
and $b$ that solves the linear equations
$$a_0 ab+a_0bP+a_1b+b_0P+b_1= 0_{1\times 4}.$$
Note that this implies
$$
b   = -[(a_0a+a_1)I_{4\times 4} +a_0 P']^{-1} 
             (b_0P+b_1I_{4\times 4})'.
$$

\bigskip
\bigskip
\subsection{Allowing for Adjustment Costs}

To allow for adjustment costs, we solve
the household maximization problem for the benchmark
model, namely \Ep{utility}
subject to \Ep{budget}, \Ep{nonneg}, and
$$
N_{t+1} k_{t+1}= [(1-\delta)k_t+x_t-\varphi(x_t/k_t)k_t]N_t
\EQN capaccnew
$$
instead of \Ep{capacc}, where
$$
\varphi(x/k) = {a\over 2} \left({x\over k}-b\right)^2.
$$
As a sensitivity check we set $b$ equal to the investment-capital
trend rate (i.e., $b=(1+g_z)(1+g_n)-1+\delta$) 
and increase $a$ from 0.  To do this, we need to
modify the dynamic first-order condition.
We replace \Ep{dynamic} by
$$\EQNalign{
 & 
   \biggl[
     {1+\tau_{xt}\over 1+\tau_{ct}}
      U_1(\hat c_t,1-l_t)-\zeta\min(\hat x_t,0)^2 
   \biggr]\left({1\over 1-\varphi'(\hat x_t/\hat k_t)}\right)\cr
     \noalign{\medskip}
 &\qquad\qquad\qquad
   = \hat \beta
     E_t\Biggl\{{U_1(\hat c_{t+1},1-l_{t+1})\over
                  1+\tau_{ct+1} }
             \bigl\{ (1-\tau_{kt+1})r_{t+1}+\delta \tau_{kt+1}\bigr\} \cr
\noalign{\medskip}
 &\qquad\qquad\qquad\qquad\qquad
    + \biggl[{1+\tau_{xt+1}\over 1+\tau_{ct+1}}
                U_1(\hat c_{t+1},1-l_{t+1})-\zeta \min(\hat x_{t+1},0)^2
         \biggr]\cr
\noalign{\smallskip}
 &\qquad\qquad\qquad\qquad\qquad\qquad
            \cdot
   \left({1\over 1-\varphi'(\hat x_{t+1}/\hat k_{t+1})}\right)\cr
\noalign{\smallskip}
 &\qquad\qquad\qquad\qquad\qquad\qquad
           \cdot \biggl\{1-\delta
           -\varphi\left({\hat x_{t+1}\over \hat k_{t+1}}\right)
           +\varphi'\left({\hat x_{t+1}\over \hat k_{t+1}}\right)
           {\hat x_{t+1}\over \hat k_{t+1}} \biggr\}
\Biggr\}.
\EQN dynamicnew\cr
}
$$

In summary, allowing for adjustment costs involves
two changes in the system of equations: replacing
\Ep{capacc} and \Ep{dynamic} 
with 
\Ep{capaccnew} and \Ep{dynamicnew}. 

For the benchmark model with adjustment costs,
we can construct 
$$
  \tilde \tau_{xt} = {1+\tau_{xt}\over 1-\varphi'(\hat x_t/\hat k_t)} -1
$$
using series on $\tau_{xt}$, $\hat x_t$, and $\hat k_t$ from
the benchmark model {\sl without} adjustment costs.
The effective investment wedge corresponding to this
rate ($1/(1+\tilde \tau_{xt})$), when fed into the benchmark model
with adjustment costs, yields almost exactly the same results 
for equilibrium output, hours, and investment.  (There is
a slight difference because the effective depreciation rates
are different for the two models.)

To do log-linear computation (as in the baseline economy) 
in the case with adjustment costs and $\tau_{ct}=\tau_{kt}=0$, 
we start with
$$\EQNalign{
 &\hat c_t + \hat g_t +(1+g_z)(1+g_n)\hat k_{t+1}-(1-\delta)\hat k_t
        +\varphi(\hat x_t/\hat k_t) \hat k_t
    = \hat y_t = \hat k_t^\theta(z_tl_t)^{1-\theta}
\EQN {foc1adj}\cr
\noalign{\medskip}
 &{\psi \hat c_t\over 1-l_t}
    = (1-\tau_{lt})(1-\theta) \hat k_t^\theta l_t^{-\theta} z_t^{1-\theta}
\EQN {foc2adj}\cr
\noalign{\medskip}
 &(1+\tau_{xt}) \hat c_t^{-\sigma} (1-l_t)^{\psi (1-\sigma)}
         /\bigl(1-\varphi'(\hat x_t/\hat k_t)\bigr)
     \cr
 &\qquad
       =  \hat \beta E_t
  \hat c_{t+1}^{-\sigma} (1-l_{t+1})^{\psi (1-\sigma)}
  \Bigl[
      \theta \hat k_{t+1}^{\theta-1} (z_{t+1}l_{t+1})^{1-\theta}
         +\bigl(1-\delta\cr
 &\qquad\qquad\qquad\qquad
            -\varphi(\hat x_{t+1}/\hat k_{t+1})
            +\varphi'(\hat x_{t+1}/\hat k_{t+1})\hat x_{t+1}/\hat k_{t+1}\bigr)
   \cr
 &\qquad\qquad\qquad\qquad
              (1+\tau_{xt+1})/
            \left(1-\varphi'(\hat x_{t+1}/\hat k_{t+1})\right)
   \Bigr].
\EQN {foc3adj}\cr
}
$$
Assuming $\varphi(\hat x/\hat k)=\varphi'(\hat x/\hat k)=0$,
the log-linearization of these equations yields the same results
as in the benchmark with the exception of the intertemporal condition:
$$ \EQNalign{
0&\approx
  (1+\tau_x)
        \bigl\{ -\psi(1-\sigma) l/(1-l) \log l_t
                -\sigma \log \hat c_t 
          +\eta (\log \hat x_t-\log\hat k_t) \bigr\} +\tau_{xt}\cr
 & \qquad 
       -E_t\Bigl\{ (1+\tau_x)
            \left\{-\psi(1-\sigma) l/(1-l) \log l_{t+1}
                   -\sigma \log \hat c_{t+1}
            \right\}\cr
 & \qquad \qquad\qquad
     +\hat\beta\bigl[r(1-\theta)
                  (\log l_{t+1} + \log z_{t+1}-\log \hat k_{t+1})\cr
 & \qquad \qquad\qquad\qquad
                  +(1+\tau_x)(1+g_z)(1+g_n)\eta
                               (\log\hat x_{t+1} - \log\hat k_{t+1})\cr
 & \qquad \qquad\qquad\qquad
                 + (1-\delta) \tau_{xt+1}\bigr]\Bigr\},
\EQN dynamicadj
\cr
}
$$
where  
$r=\theta \hat y/\hat k$
and $\eta=\varphi''(\hat x/\hat k)(\hat x/\hat k)=ab$.
          
As before, we guess a solution of the form \Ep{formofsolution}
and set $\gamma$'s so that the dynamic residual \Ep{dynamicadj} is 
exactly 0.
We start by deriving $\gamma_k$.  To do this, 
we need to write out the coefficients on $\hat k_{t+2}$, 
$\hat k_{t+1}$, and $\hat k_t$ in the dynamic first-order condition
\Ep{dynamicadj} making use
of the $\phi$'s from the static first-order conditions. 
For now, we 
can ignore the expectations operator.  We get the following:
$$\EQNalign{
  0 = &\bigl[\hat\beta r(1-\theta)\phi_{lk'}
                 -(1+\tau_x)
                    \{\psi(1-\sigma)l/(1-l)\phi_{lk'}
                      +\sigma\phi_{ck'}
                      -\hat\beta (1+g_n)(1+g_z)
                         \eta\phi_{xk'} \} \bigr] \hat k_{t+2}\cr
      +&\bigl[\hat\beta r(1-\theta)(\phi_{lk}-1)
               -(1+\tau_x)
                     \{\psi(1-\sigma)l/(1-l)(\phi_{lk}-\phi_{lk'})
                       + \sigma(\phi_{ck}-\phi_{ck'}) \cr
       &\qquad\qquad\qquad\qquad
                       -\hat\beta (1+g_n)(1+g_z) \eta (\phi_{xk}-1)
                        +\eta\phi_{xk'} \}
                  \bigr] \hat k_{t+1}\cr
      +&\bigl[(1+\tau_x)
               \{\psi(1-\sigma)l/(1-l)\phi_{lk}
                +\sigma\phi_{ck} -
                     \eta (\phi_{xk}-1)\} \bigr] \hat k_t\cr
      +&{\rm \ all\ other\ terms}
\EQN dynamicpieceadj
}
$$
or, more succinctly, rewrite \Ep{dynamicpieceadj} as
$ (a+ bL + cL^2) \hat k_{t+2} =$ other terms,
where $\gamma_k$ is the root of the quadratic 
inside the unit circle.  Note that $\gamma_k$ does not depend
on the other unknown $\gamma$'s.

Given $\gamma_k$, we can solve a linear system
for the other $\gamma$'s.  At this point,
we do not ignore the expectations operator:
$$ \eqalign{
0&\approx
  (1+\tau_x)
        \bigl\{-\psi(1-\sigma) l/(1-l)
              [(\phi_{lz}+\phi_{lk'}\gamma_z)\log z_t
               +(\phi_{ll}+\phi_{lk'}\gamma_l)\tau_{lt}\cr
 & \qquad\qquad\qquad \qquad\qquad\qquad\qquad
               +\phi_{lk'}\gamma_x\tau_{xt} 
               +(\phi_{lg}+\phi_{lk'}\gamma_g)\log\hat g_t]\cr
 & \qquad\qquad\qquad \qquad
               -\sigma 
              [(\phi_{cz}+\phi_{ck'}\gamma_z)\log z_t
               +(\phi_{cl}+\phi_{ck'}\gamma_l) \tau_{lt}\cr
 & \qquad\qquad\qquad \qquad\qquad\qquad
               +\phi_{ck'}\gamma_x \tau_{xt}
               +(\phi_{cg}+\phi_{ck'}\gamma_l) \log\hat g_t
               ]\cr
 & \qquad\qquad\qquad \qquad\qquad\qquad
             +\eta \phi_{xk'}[\gamma_z \log z_t + \gamma_l \tau_{lt}
                                + \gamma_x \tau_{xt} + \gamma_g\log\hat g_t]
             \bigr\}+\tau_{xt}\cr
}
$$
$$\eqalign{
 & \qquad 
       -E_t(1+\tau_x)
            \bigl\{-\psi(1-\sigma) l/(1-l) 
              [(\phi_{lk}+\phi_{lk'}\gamma_k)(\gamma_z\log z_t
                                              +\gamma_l\tau_{lt}
                                              +\gamma_x\tau_{xt}
                                              +\gamma_g\log\hat g_t)\cr
 & \qquad\qquad\qquad \qquad\qquad\qquad\qquad \qquad
              + (\phi_{lz}+\phi_{lk'}\gamma_z)\log z_{t+1}\cr
 & \qquad\qquad\qquad \qquad\qquad\qquad\qquad \qquad
              + (\phi_{ll}+\phi_{lk'}\gamma_l)\tau_{lt+1}\cr
 & \qquad\qquad\qquad \qquad\qquad\qquad\qquad \qquad
              + \phi_{lk'}\gamma_x\tau_{xt+1}\cr
 & \qquad\qquad\qquad \qquad\qquad\qquad\qquad \qquad
              + (\phi_{lg}+\phi_{lk'}\gamma_g)\log\hat g_{t+1}]\cr
}
$$
$$\eqalign{
 & \qquad\qquad\qquad\qquad\qquad \qquad
               -\sigma 
              [(\phi_{ck}+\phi_{ck'}\gamma_k)(\gamma_z\log z_t
                                              +\gamma_l\tau_{lt}
                                              +\gamma_x\tau_{xt}
                                              +\gamma_g\log\hat g_t)\cr
 & \qquad\qquad\qquad\qquad\qquad \qquad\qquad
               +(\phi_{cz}+\phi_{ck'}\gamma_z)\log z_{t+1}\cr
 & \qquad\qquad\qquad\qquad\qquad \qquad\qquad
               +(\phi_{cl}+\phi_{ck'}\gamma_l)\tau_{lt+1}\cr
 & \qquad\qquad\qquad\qquad\qquad \qquad\qquad
               +\phi_{ck'}\gamma_x\tau_{xt+1}\cr
 & \qquad\qquad\qquad\qquad\qquad \qquad\qquad
               +(\phi_{cg}+\phi_{ck'}\gamma_g)\log\hat g_{t+1}]\bigr\}\cr
}
$$
$$\eqalign{
 & \qquad \qquad\quad
     -\hat\beta r(1-\theta) E_t \bigl[(\phi_{lk}+\phi_{lk'}\gamma_k-1)
                    (\gamma_z\log z_t + \gamma_l\tau_{lt}
                               +\gamma_x\tau_{xt}
                               +\gamma_g\log\hat g_t)\cr
 & \qquad\qquad\qquad\qquad\qquad \qquad\qquad
                    +(\phi_{lz}+\phi_{lk'}\gamma_z+1)\log z_{t+1}\cr
 & \qquad\qquad\qquad\qquad\qquad \qquad\qquad
                    +(\phi_{ll}+\phi_{lk'}\gamma_l)\tau_{lt+1}\cr
 & \qquad\qquad\qquad\qquad\qquad \qquad\qquad
                    +\phi_{lk'}\gamma_x\tau_{xt+1}\cr
 & \qquad\qquad\qquad\qquad\qquad \qquad\qquad
                    +(\phi_{lg}+\phi_{lk'}\gamma_g)\log\hat g_{t+1}\bigr]\cr
}$$
$$\eqalign{
 & \qquad \qquad\quad
   -\hat\beta (1+\tau_x)(1+g_n)(1+g_z)\eta 
           E_t\bigl[(\phi_{xk}+\phi_{xk'}\gamma_k-1)
                    (\gamma_z\log z_t + \gamma_l\tau_{lt}
                               +\gamma_x\tau_{xt}
                               +\gamma_g\log\hat g_t)\cr
 & \qquad\qquad\qquad\qquad\qquad \qquad\qquad
                    +\phi_{xk'}(\gamma_z\log z_{t+1}
                              + \gamma_l\tau_{lt+1}
                               +\gamma_x\tau_{xt+1}
                               +\gamma_g\log\hat g_{t+1})\cr
 & \qquad\qquad\qquad
     -\hat \beta (1-\delta) E_t \tau_{xt+1},\cr
}
$$
which reduces to a system like \Ep{lqsystem}.

\subsection{An Alternative Investment Wedge}

In a comment on our paper, Christiano and Davis (2006)
note that our findings may be sensitive to the particular
choice of the intertemporal wedge.  In theory the choice
of the intertemporal wedge should not
matter in the sense that it just implies a (slightly)
different map between the detailed economy and the
prototype economy.  However, in practice, stochastic
processes for the wedges are estimated, and therefore
we want to make sure that our substantive
findings are not affected by the choice.
Here we describe Christiano and Davis' (2006) 
alternative investment wedge,
and later we demonstrate that our findings are 
not sensitive to this alternative.\note{We also explain later 
why the methodology used by Christiano and 
Davis (2006)---which is not the methodology used in 
the final version of our paper---can lead to
a different conclusion.}

Christiano and Davis (2006) assume that the intertemporal
wedge is $\tau_k$ and that the dynamic first-order condition
(with adjustment costs) is given by
$$\EQNalign{
 & U_1(\hat c_t,1-l_t)/(1-\varphi'(\hat x_t/\hat k_t))\cr
\noalign{\medskip}
 &\qquad = \hat \beta
     E_t\Biggl\{U_1(\hat c_{t+1},1-l_{t+1})
(1-\tau_{kt+1}) \Bigl[r_{t+1}
   +
   {1\over 1-\varphi'(\hat x_{t+1}/\hat k_{t+1})}\cr
\noalign{\medskip}
 &\qquad \qquad\qquad\qquad\qquad \qquad\qquad\qquad
           \cdot 
          \biggl\{1-\delta
           -\varphi\left({\hat x_{t+1}\over \hat k_{t+1}}\right)
           +\varphi'\left({\hat x_{t+1}\over \hat k_{t+1}}\right)
           {\hat x_{t+1}\over \hat k_{t+1}} \biggr\} \Bigr]\Biggr\}.\cr
}
$$
Here, $\tau_k$ resembles a tax on the gross return to capital.
In Appendix C, we compare our predictions based on
the $\tau_k$ wedge with those based on the $\tau_x$ wedge.\note{%
We report results only for the postwar period 
applying a log-linear approximation method, so we have dropped
the penalty functions needed for our nonlinear solution method.}

\bigskip
\section{MLE Estimation}

We now describe the general method we use to 
estimate the processes governing
the four exogenous variables in $s_t$ with the data
described above.  

\bigskip
\subsection{State-Space Form}
$$\eqalign{
     X_{t+1} & = A  X_t + B \epsilon_{t+1} \cr
     Y_t     & = C  X_t + \omega_t\cr
     \omega_t & = D \omega_{t-1} + \eta_t,\cr
}
$$
\item{ }
where $X_t=[\log \hat k_t, \log z_t, \tau_{lt}, \tau_{xt}, \log\hat g_t, 1]'$,
$Y_t=[\log\hat y_t,\log \hat x_t,\log l_t, \log\hat g_t]$, and
$$\eqalign{
   A &= \left[\matrix{ \gamma_k  & 
                        \matrix{\gamma_z & \gamma_l & \gamma_x & \gamma_g} 
                           & \gamma_0 \cr
                 0_{4\times 1}  & P       &  P_0 \cr
                 0              & 0_{1\times 4} & 1}\right]\cr
\noalign{\bigskip}
   B &= \left[\matrix{  0_{1\times 4} \cr
                       Q \cr
                       0 }\right]\cr
\noalign{\bigskip}
   C &= \left[\matrix{
     \phi_{yk} & \phi_{yz} & \phi_{yl} & 0 & \phi_{yg} & \phi_{y0}\cr
     \phi_{xk} & 0 & 0 & 0 & 0 & \phi_{x0}\cr
     \phi_{lk} & \phi_{lz} & \phi_{ll} & 0 & \phi_{lg} & \phi_{l0}\cr
     0 & 0 & 0 & 0 & 1 & 0\cr
    }\right]
   +\left[\matrix{\phi_{yk'}\cr \phi_{xk'}\cr \phi_{lk'}\cr 0\cr}\right]
   \left[\matrix{\gamma_k & \gamma_z & \gamma_l & \gamma_x & \gamma_g & 0\cr}
    \right]
}
$$
\item{}
and elements of $D$ are the parameters governing serial correlation of
the measurement error.  Assume that $E \eta_t\eta_t'=R$, 
$E \epsilon_{t}\eta_s'=0$ for all periods $t$ and $s$.
Define $\bar Y_t\equiv Y_{t+1}-DY_t$. Then we can rewrite the 
system as
$$\eqalign{
     X_{t+1} & = A  X_t + B \epsilon_{t+1} \cr
     \bar Y_t  & = \bar C  X_t + CB\epsilon_{t+1}+\eta_{t+1}.\cr
}
$$

\subsection{Log-Likelihood Function}
$$L(\Theta) = \sum_{t=0}^{T-1} \left\{ \log |\Omega_t| + 
              {\rm trace}(\Omega_t^{-1}u_tu_t')
             -\log|\partial f(Z_t,\Theta)/\partial Z_t|\right\},
\EQN likelihood
$$
where the parameters to be estimated are stacked in vector
$\Theta$, the innovation vector is $u_t$, and its covariance
is $\Omega_t$.  The last term in \Ep{likelihood} is nonzero
if the elements of 
$Y$ are not the raw series but depend on the raw
series $Z$ plus the parameter vector.  For example, if 
we estimate $g_z$ and use per-capita values as our raw data,
then $Z$ is per-capita data and $Y$ is detrended per-capita data.   

The innovation vector $u_t$ and its covariance $\Omega_t$
are defined as follows:
$$\eqalign{
   u_t &= \bar Y_t - \hat E[\bar Y_t| 
              \bar Y_{t-1},\bar Y_{t-2},\ldots, \bar Y_0,\hat X_0]\cr
       &= Y_{t+1} - \hat E[Y_{t+1}| Y_t,Y_{t-1},\ldots, Y_0,\hat X_0]\cr
       &= Y_{t+1} - D Y_t - \bar C \hat X_t \cr
\noalign{\bigskip}
  \Omega_t & = Eu_tu_t'= \bar C \Sigma_t\bar C'+R+CBB'C',\cr
}
$$
which in turn depends on the predicted state $\hat X_t$:
$$
  \hat X_t = \hat E[X_t| Y_t,Y_t,\ldots, Y_0,\hat X_0].
$$
The predicted state evolves according to 
$$
  \hat X_{t+1} = A \hat X_t + K_t u_t, 
$$
where $K_t$ is the Kalman gain,
$$\eqalign{
  K_t &= (BB'C'+A\Sigma_t\bar C')\Omega_t^{-1}\cr
  \Sigma_{t+1} & = A\Sigma_t A'+BB'-(BB'C'+A\Sigma_t\bar C')\Omega_t^{-1}
           (\bar C\Sigma_t A'+CBB')\cr
}
$$
with state covariance $\Sigma_t$.

For the results in the paper, we fixed parameters
of preferences, production, and growth and estimated the
processes for the wedges.  The parameters that
were fixed were $\psi=2.24$, $\sigma=1$, $\beta=0.9722$, 
$\theta=0.35$,
$\delta=0.0464$, $g_n=1.5$\%, and $g_z=1.6$\%.
We also set the measurement errors equal to zero ($D=R=0_{4\times 4}$)
in all periods and all numerical experiments.\note{This choice makes no
difference for our results.}
The parameters that were estimated were elements
of $P_0$, $P$, and $Q$.   

The parameter choices were based on time series during the 
pre-World War II period.  There have been some modest 
changes in growth rates, depreciation rates, and capital shares
in the post-World War II period.  Since we separately estimate
the means of the wedges in the pre- and postwar periods, our
main results are not affected.  One advantage of keeping the
utility and technology parameters fixed is that we can back out
wedges for the entire century.  (These are shown later in Figures
A1 and A2.)

\bigskip
\section{Decomposing Macro Aggregates}

Here, we describe the details of our implementing
our accounting procedure for the Great Depression
and postwar periods.   They differ slightly because
the equilibria are computed differently in the two
periods: we use a nonlinear computational routine
for the Great Depression and a log-linear computational
routine for the postwar period.

\subsection{Great Depression Period}

Because shocks are large during the Great Depression 
period, we need to compute equilibria nonlinearly.
We use a large (i.e., 459 state) Markov chain to approximate
the process for the shocks to avoid
having to compute a five-dimensional continuous-state
model.  

Specifically, using estimates for the matrices underlying the 
stochastic process on wedges, namely $P_0$, $P$,
and $Q$, we construct a Markov transition matrix
using a methodology that is similar to---but not the
same as---Tauchen (1986).
The difference is that Tauchen rewrites the problem
to have a diagonal covariance matrix on the disturbances
and we assume a dense covariance matrix.  We do this
because it allows us more flexibility in choosing the
grid over our state variables.

Given a Markov chain, we want to find the realization
of the wedges, one consistent with the estimated
process, that implies exact agreement between model
simulations and observations.
For $\log z_t$, we use the specified production function
above and observations on output, labor input, and the
capital stock (accumulated via perpetual inventory with
the observed investment series).
For $\tau_{lt}$, we use the static first-order
condition \Ep{foc2} along with series for
capital, labor, consumption, and $z_t$.
For $\log g_t$, we have a direct measure.

We cannot infer a realization for $\tau_{xt}$ directly
from static first-order conditions.   Instead, we
find the realization of $\tau_{xt}$ that gives an
exact match of the model simulation and U.S.~observations
over a specified period (e.g., the Great Depression
or the 1982 recession).
We start by making a guess for the series $\tau_{xt}$ 
for the period in which we are interested.  For example, when
we examine the period of the Great Depression, we 
make a guess about the realization of $\tau_{xt}$
from 1929 through 1939.    A good starting point for
this guess can be derived from the log-linear decision
rule on investment, which is a function of the capital
stock and the other wedges.    This amounts to solving
one equation in one unknown ($\tau_x$).

With a guess for the realization of $\tau_{xt}$ 
and realizations for all of the other wedges, we simulate
the model as follows.   We take our Markov chain for
the wedges (found using the variation on Tauchen's method)
and append it by $N$ states, where $N$ 
is the number of years in the period we want to study.
Suppose that the 
original chain had $M$ states.  The appended chain has $M+N$.
In the case
of the Great Depression, for example, we have $N=11$,
which is the length of the period 1929--1939.
The first of the appended states corresponds to the
first year; the first year is 1929 in the case of the Great Depression.
We have values for $z_{1929}$, $\tau_{l1929}$, $\hat g_{1929}$. 
We have a guess for $\tau_{x1929}$.  
For each of the $N$ additional states, we use our version
of Tauchen's method to compute the transition problems from 
this state to the $M$ original states.
Because we want a sequence for $\tau_{xt}$ during the Great
Depression that generates an exact match between data and 
model for all wedges on, we have to iterate (reguessing
the sequence and updating the Markov chain transition)
until there is an exact match.
The derived wedges are shown in Figure 1 of the paper.

Once we have this realization for all of the wedges---one that
can exactly generate the data---we turn off one or more
of the wedges by setting them at their initial levels (say, 1929
in the case of the Depression).    For example, to
see how important $\log z$ is in the Great Depression, 
we can set $\tau_{x}$, $\tau_{l}$, and $\hat g$ (in 
all states) equal to their 1929 levels.  We hold fixed
the underlying stochastic process.  This means that
we hold the Markov chain transition probabilities fixed
for each new simulation.  
The results of the one-wedge-alone or one-wedge-off
experiments are shown in Figures 2--4.

\subsection{Postwar Period}

In the postwar period, there is no need to approximate
the stochastic process for the wedges using a Markov
Chain because the shocks are much smaller. In this case,
a log-linear approximation of the five-dimensional,
continuous state model works well.  

As in the Great
Depression, we fix the stochastic process $s^t$
when considering the marginal effect of one wedge. 
To accomplish that, we set each of the wedges equal
to the sum of a constant times an indicator function
plus the corresponding state in $s$ times one minus
the indicator function; for example,
$$
  \tau_{lt} = \chi \bar\tau_l + (1-\chi) s_{2t},
$$
where $s_{2t}$ is the second element of the 
state vector $s_t$ and $\chi$ is equal to 0 or
1 depending on whether variations in the
wedge $\tau_{lt}$ are being analyzed 
or not, respectively.\note{In earlier versions of
the paper we set $\tau_{lt}=s_{2t}$ and 
set it equal to a constant when its effects were
not being analyzed.  The quantitative impact was
very small for our prototype model with a $\tau_x$
investment wedge.  They are not small for Christiano
and Davis (2006), who prefer to use a $\tau_k$ wedge and 
to set adjustment
costs very high.   In the published version of our paper, 
we described our current procedure which is
consistent with our propositions that separate 
the direct effects of fluctuations in the wedges with 
indirect effects due to forecasting fluctuations 
in other wedges.  Later, we show that if Christiano
and Davis (2006) were to apply the procedure 
that is consistent with our
propositions, their predictions would line up 
almost exactly with ours. }

\bigskip
\section{A Model with Capacity Utilization}

In the cases where we allow for capacity utilization
(Figures 9--12 in the paper), we
use the following functional form for technology:
$$F(k,zl) = k^\theta zl$$
and redo the steps outlined above.

In the log-linearized solution, we have the following new
first-order conditions:
$$\eqalign{
 & \hat c_t +\hat g_t+(1+g_n)(1+g_z)\hat k_{t+1}-(1-\delta)\hat k_t
   = \hat k_t^\theta z_tl_t\cr
 &\psi \hat c_t
    = (1-\tau_{lt})\hat k_t^\theta z_t (1-l_t) \cr
 &(1+\tau_{xt}) \hat c_t^{-\sigma}
              (1-l_t)^{\psi (1-\sigma)}
     \cr
 &\qquad\qquad
       = \hat \beta E_t\hat c_{t+1}^{-\sigma}
                (1-l_{t+1})^{\psi (1-\sigma)}\cr
 &\qquad\qquad\qquad\qquad
          [\theta \hat k_{t+1}^{\theta-1} z_{t+1}l_{t+1}
              +(1-\delta)(1+\tau_{xt+1})].\cr
}
$$
The new steady state in this case is $\hat k$ and $l$ that solve
$$\eqalign{
  & \psi [\hat k^\theta zl
             -(1+g_n)(1+g_z)\hat k+(1-\delta)\hat k-\hat g]
    = (1-\tau_l)\hat k^\theta z (1-l) \cr
 &(1+\tau_x) = \hat \beta [\theta \hat k^{\theta-1} zl
              +(1-\delta)(1+\tau_x)].\cr
}
$$
Solving this system of equations is 
like solving a problem of the form
$ \xi_1 \hat k - \xi_2 = \xi_3 \hat k^\theta$,
where 
$$\eqalign{
\xi_1&= (\psi+1-\tau_l)\xi-\psi[(1+g_n)(1+g_z)-1+\delta]\cr
\xi_2&=\psi \hat g\cr
\xi_3&=(1-\tau_l)z\cr
\xi\  & =(1+\tau_x)[1-\hat \beta(1-\delta)]/(\hat\beta\theta).\cr
}
$$
We compute $\hat k$ with 
a simple Newton algorithm starting with
$\hat k= (\xi_3/\xi_1)^{1/(1-\theta)},$
which is the exact solution if $\hat g=0$.
The solution for $\hat k$ 
in the $\hat g > 0$ case will be higher.
Once we have $\hat k$,  we have
$l=\xi \hat k^{1-\theta}/z$.

The new log-linearized first-order conditions (ignoring
constants) imply the following for detrended consumption:
$$\eqalign{
  \hat c_t &\approx \hat c \log \hat c_t\cr
           &\approx \hat k^{\theta} zl
           [\theta \log \hat k_t + \log z_t+\log l_t]\cr
          &\quad
          -(1+g_z)(1+g_n)\hat k \log \hat k_{t+1} 
          +(1-\delta)\hat k \log \hat k_t-\hat g \log\hat g_t. \cr
}
$$
As before, the labor input is derived from the static
first-order condition \Ep{foc2}
$$\eqalign{
    0 &\approx \psi\bigl\{\hat k^{\theta} zl
           [\theta \log \hat k_t + \log z_t+\log l_t]\cr
          &\qquad\quad
          -(1+g_z)(1+g_n)\hat k \log \hat k_{t+1} 
          +(1-\delta)\hat k \log \hat k_t-\hat g\log \hat g_t\bigr\} \cr
          &\qquad\qquad\quad 
                 + (1-\tau_l)\hat k^\theta 
                          z(1-l)
                    \bigl\{1/(1-\tau_l)\,\tau_{lt}\cr
          &\qquad\qquad\qquad\quad 
                               -\theta \log \hat k_t
                                  - \log z_t
                                  +l/(1-l) \log l_t \bigr\},\cr
}
$$
which we write succinctly as
$$
\log l_t = \phi_{lk}\log \hat k_t
            +\phi_{lz}\log z_t
            +\phi_{ll}\tau_{lt}
            +\phi_{lg}\log \hat g_t
            +\phi_{lk'}\log \hat k_{t+1}.
$$
Using this equation for $\log l$, we use the other static
first-order conditions to derive output, investment,
and consumption as follows:
$$\eqalign{
\log \hat y_t &= \phi_{yk}\log \hat k_t
            +\phi_{yz}\log z_t
            +\phi_{yl}\tau_{lt}
            +\phi_{yk'}\log \hat k_{t+1} \cr
      &= 
              (\theta+\phi_{lk})\log \hat k_t
                        +(1+\phi_{lz}) \log z_t \cr
      &\qquad\qquad\qquad\qquad
                        +\phi_{ll}\tau_{lt}
                                 +\phi_{lk'}\log \hat k_{t+1}\cr
      \log \hat x_t &=
                    (1+g_z)(1+g_n) \hat k/\hat x \log \hat k_{t+1}
                    -(1-\delta) \hat k/\hat x \log \hat k_t\cr
      \log \hat c_t &= \phi_{ck}\log \hat k_t
            +\phi_{cz}\log z_t
            +\phi_{cl}\tau_{lt}
            +\phi_{cg}\log \hat g_t
            +\phi_{ck'}\log \hat k_{t+1} \cr
                    &= [\hat y\log y_t
                      -\hat x\log x_t -\hat g\log \hat g_t]/\hat c.\cr
}
$$

The new dynamic first-order condition implies the following
for capital:
$$ \eqalign{
0&\approx
  (1+\tau_x)\hat c^{-\sigma} (1-l)^{\psi(1-\sigma)} 
        \left\{-\psi(1-\sigma) l/(1-l) \log l_t
                  -\sigma \log \hat c_t \right\} \cr
 & \qquad \qquad\qquad
   +\hat c^{-\sigma} (1-l)^{\psi(1-\sigma)} \tau_{xt} \cr
 & \qquad 
   -\hat\beta 
        E_t\bigl\{ \bigl[\theta \hat k^{\theta-1}zl
           +(1-\delta)(1+\tau_x)\bigr]\cr
 & \qquad \qquad\qquad
       \cdot\bigl[\hat c^{-\sigma} (1-l)^{\psi(1-\sigma)}
            \left\{-\psi(1-\sigma) l/(1-l) \log l_{t+1}
                    -\sigma \log \hat c_{t+1} \right\}
       \bigr]\cr
 & \qquad \qquad\qquad
        +\hat c^{-\sigma} (1-l)^{\psi(1-\sigma)}
             \bigl[\theta \hat k^{\theta-1}zl
        \cr
 & \qquad \qquad\qquad\qquad\qquad
       \cdot(\log l_{t+1} + \log z_{t+1}-(1-\theta)\log \hat k_{t+1})
                  +(1-\delta)\tau_{xt+1}\bigl]\bigr\},\cr
}
$$
which simplifies to
$$ \eqalign{
0&\approx
  (1+\tau_x)
        \left\{ -\psi(1-\sigma) l/(1-l) \log l_t
                -\sigma \log \hat c_t \right\} 
   +\tau_{xt} \cr
 & \qquad 
       -E_t\bigl\{ (1+\tau_x)
            \left\{-\psi(1-\sigma) l/(1-l) \log l_{t+1}
                   -\sigma \log \hat c_{t+1}
            \right\}\cr
 & \qquad \quad
     +\hat\beta\bigl[r
                  (\log l_{t+1} + \log z_{t+1}-(1-\theta)\log \hat k_{t+1})
                  +(1-\delta)\tau_{xt+1}\bigr]\bigr\},\cr
}
$$
where $r=\theta \hat y/\hat k$.
          
The form of the solution is given by \Ep{formofsolution}.
To compute $\gamma_k$,
we need to write out the coefficients on $\hat k_{t+2}$, 
$\hat k_{t+1}$, and $\hat k_t$ in the dynamic first-order
condition making use
of the $\phi$'s from the static first-order condition.  For now, we 
can ignore the expectations operator.  We get the following:
$$ \eqalign{
0&\approx
  -(1+\tau_x)
        \bigl\{ -\psi(1-\sigma) l/(1-l)
              [\phi_{lk}\log \hat k_t+\phi_{lk'} \log \hat k_{t+1}] 
   \cr
 & \qquad\qquad\qquad \qquad
               -\sigma 
              [\phi_{ck}\log \hat k_t+\phi_{ck'} \log \hat k_{t+1}] 
         \bigr\} \cr
 & \quad 
       +E_t(1+\tau_x)
            \bigl\{-\psi(1-\sigma) l/(1-l) 
              [\phi_{lk}\log \hat k_{t+1}+\phi_{lk'} \log \hat k_{t+2}] \cr
 & \qquad\qquad\qquad\qquad\qquad \qquad
               -\sigma 
              [\phi_{ck}\log \hat k_{t+1}+\phi_{ck'} \log \hat k_{t+2}] 
            \bigr\}\cr
 & \qquad \qquad\quad
     +\hat\beta E_t \bigl[r
                  (\phi_{lk}\log \hat k_{t+1}+\phi_{lk'}\hat k_{t+2}
                     -(1-\theta)\log \hat k_{t+1})
              \bigr],\cr
}
$$
which simplifies to
$$\eqalign{
  0 = &[\hat\beta r\phi_{lk'}
                 -(1+\tau_x)\psi(1-\sigma)l/(1-l)\phi_{lk'}
         -(1+\tau_x)\sigma\phi_{ck'}] \hat k_{t+2}\cr
      +&[\hat\beta r(\phi_{lk}-1+\theta)
               -(1+\tau_x)\psi(1-\sigma)l/(1-l)\phi_{lk}
               -(1+\tau_x)\sigma\phi_{ck}
       \cr
       &\qquad\qquad\qquad\qquad
               +(1+\tau_x)\psi(1-\sigma)l/(1-l)\phi_{lk'}
               +(1+\tau_x)\sigma\phi_{ck'}
       ] \hat k_{t+1}\cr
      -&[-(1+\tau_x)\psi(1-\sigma)l/(1-l)\phi_{lk}
               -(1+\tau_x)\sigma\phi_{ck}
       ] \hat k_t\cr
      +&{\rm\ all\ other\ terms.}
}
$$

Given $\gamma_k$, we solve a linear system
for the other $\gamma$'s.   In particular, we have
$$ \eqalign{
0&\approx
  (1+\tau_x)
        \bigl\{-\psi(1-\sigma) l/(1-l)
              [(\phi_{lz}+\phi_{lk'}\gamma_z)\log z_t\cr
 & \qquad\qquad\qquad \qquad\qquad\qquad\qquad
               +(\phi_{ll}+\phi_{lk'}\gamma_l)\tau_{lt}\cr
 & \qquad\qquad\qquad \qquad\qquad\qquad\qquad
               +\phi_{lk'}\gamma_x\tau_{xt} 
               +(\phi_{lg}+\phi_{lk'}\gamma_g)\log\hat g_t\cr
 & \qquad\qquad\qquad \qquad
               -\sigma 
              [(\phi_{cz}+\phi_{ck'}\gamma_z)\log z_t
               +(\phi_{cl}+\phi_{ck'}\gamma_l) \tau_{lt}\cr
 & \qquad\qquad\qquad \qquad\qquad\qquad
               +\phi_{ck'}\gamma_x \tau_{xt}
               +(\phi_{cg}+\phi_{ck'}\gamma_l) \log\hat g_t
               ]\bigr\} +\tau_{xt}\cr
}
$$
$$\eqalign{
 & \qquad 
       -E_t(1+\tau_x)
            \bigl\{-\psi(1-\sigma) l/(1-l) 
              [(\phi_{lk}+\phi_{lk'}\gamma_k)(\gamma_z\log z_t
                                              +\gamma_l\tau_{lt}
                                              +\gamma_x\tau_{xt})\cr
 & \qquad\qquad\qquad \qquad\qquad\qquad\qquad \qquad
              + (\phi_{lz}+\phi_{lk'}\gamma_z)\log z_{t+1}\cr
 & \qquad\qquad\qquad \qquad\qquad\qquad\qquad \qquad
              + (\phi_{ll}+\phi_{lk'}\gamma_l)\tau_{lt+1}\cr
 & \qquad\qquad\qquad \qquad\qquad\qquad\qquad \qquad
              + \phi_{lk'}\gamma_x\tau_{xt+1}\cr
 & \qquad\qquad\qquad \qquad\qquad\qquad\qquad \qquad
              + (\phi_{lg}+\phi_{lk'}\gamma_g)\log\hat g_{t+1}]\cr
}
$$
$$\eqalign{
 & \qquad\qquad\qquad\qquad\qquad \qquad
               -\sigma 
              [(\phi_{ck}+\phi_{ck'}\gamma_k)(\gamma_z\log z_t
                                              +\gamma_l\tau_{lt}
                                              +\gamma_x\tau_{xt}
                                              +\gamma_g\log\hat g_t)\cr
 & \qquad\qquad\qquad\qquad\qquad \qquad\qquad
               +(\phi_{cz}+\phi_{ck'}\gamma_z)\log z_{t+1}\cr
 & \qquad\qquad\qquad\qquad\qquad \qquad\qquad
               +(\phi_{cl}+\phi_{ck'}\gamma_l)\tau_{lt+1}\cr
 & \qquad\qquad\qquad\qquad\qquad \qquad\qquad
               +\phi_{ck'}\gamma_x\tau_{xt+1}\cr
 & \qquad\qquad\qquad\qquad\qquad \qquad\qquad
               +(\phi_{cg}+\phi_{ck'}\gamma_g)\log\hat g_{t+1}]\bigr\}\cr
}
$$
$$\eqalign{
 & \qquad \qquad\quad
     -\hat\beta r E_t \bigl[(\phi_{lk}+\phi_{lk'}\gamma_k-1+\theta)
                    (\gamma_z\log z_t + \gamma_l\tau_{lt}+\gamma_x\tau_{xt})\cr
 & \qquad\qquad\qquad\qquad\qquad \qquad\qquad
                    +(\phi_{lz}+\phi_{lk'}\gamma_z+1)\log z_{t+1}\cr
 & \qquad\qquad\qquad\qquad\qquad \qquad\qquad
                    +(\phi_{ll}+\phi_{lk'}\gamma_l)\tau_{lt+1}\cr
 & \qquad\qquad\qquad\qquad\qquad \qquad\qquad
                    +\phi_{lk'}\gamma_x\tau_{xt+1}\cr
 & \qquad\qquad\qquad\qquad\qquad \qquad\qquad
                    +(\phi_{lg}+\phi_{lk'}\gamma_g)\log\hat g_{t+1}\bigr]\cr
 & \qquad\qquad\qquad
     -\hat \beta (1-\delta) E_t \tau_{xt+1},\cr
}
$$
and after this, the procedure is the same as for the benchmark
model.

\Appendix{B}{Data and Sources}

\noindent
In this Appendix, we describe our data and their
sources.  Early drafts of the paper reported 
results for the annual data described in Section 1.
We subsequently redid the exercises for the postwar
period using quarterly data (which can be updated
by future users of the codes).  The Bureau of 
Economic Analysis (BEA) has done comprehensive revisions
of the national accounts and reenumerated many of the 
standard national income and product account (NIPA) 
tables.  Thus, the NIPA tables with
annual data do not correspond to the current
naming scheme of the BEA.  We do have all of
the original tables that we used in this project,
along with documentation.

\section{U.S.~Historical Annual Data Measures and Sources}

Here, we provide a list of the variables in the 
model and their data analogs.  Because we work
with data going back to 1900, the main output series 
for the annual data is gross national product.

\subsection{Measures}

\item{\bul} Per-capita output ($y$)
\itemitem{    } \ \ GNP 
\itemitem{    } \ \ $-$ Sales tax 
\itemitem{    } \ \ $-$ Military compensation
\itemitem{    } \ \ $+$ Services from consumer durables (with return = 4\% )
\itemitem{    } \ \ + Depreciation from consumer durables
\bigskip
\item{\bul} Per-capita investment ($x$)
\itemitem{    } \ \ Gross private fixed investment
\itemitem{    } \ \ + Private inventories 
\itemitem{    } \ \ + Government gross investment
\itemitem{    } \ \ + Net factor payments (GNP-GDP)
\itemitem{    } \ \ + Personal consumption expenditures on durables
\itemitem{    } \ \ $-$ Sales tax $\times$ share of durables in PCE
\itemitem{    } \ \ All deflated by the GNP deflator and Population over 16
\bigskip
\item{\bul} Per-capita government ($g$)
\itemitem{    } \ \ Government consumption
\itemitem{    } \ \ $-$ Military compensation
\itemitem{    } \ \ +   Military equipment
\itemitem{    } \ \ $-$ 1/2  Military facilities
\itemitem{    } \ \ + Net exports of goods and services
\itemitem{    } \ \ All deflated by the GNP deflator and Population over 16
\bigskip
\item{\bul} Per-capita labor input ($l$)
\itemitem{    } \ \ (Civilian annual manhours / Population over 16)
                       / (50 weeks $\times$ 100 hours)

\subsection{Sources}

The specific sources of the data listed above are as follows:
\medskip
\itemitem{\cir} National accounts, pre-1929
\itemitem{    } \hskip .5truein Kendrick (1961), Table A-IIb (all mil.~\$)
\itemitem{    } \hskip 1.truein Total consumption expenditures
\itemitem{    } \hskip 1.truein New construction and equipment
\itemitem{    } \hskip 1.truein Change in business inventories
\itemitem{    } \hskip 1.truein Net foreign investment
\itemitem{    } \hskip 1.truein Government purchases of goods and services
\itemitem{    } \hskip 1.truein GNP (commerce concept)
\medskip
\itemitem{    } \hskip .5truein Kendrick (1961), Table A-IIa (in mil.~1929 \$)
\itemitem{    } \hskip 1.truein GNP 
\itemitem{    } \hskip 1.truein Change in business inventories
\bigskip
\itemitem{\cir} National accounts, post-1929
\itemitem{    } \hskip .5truein www.bea.gov, old NIPA Table 1.1 (in bil.~\$)
\itemitem{    } \hskip 1.truein Gross domestic product (GDP)
\itemitem{    } \hskip 1.truein Personal consumption expenditures (PCE)
\itemitem{    } \hskip 1.5truein PCE durable goods
\itemitem{    } \hskip 1.5truein PCE nondurable goods
\itemitem{    } \hskip 1.5truein PCE services
\itemitem{    } \hskip 1.truein Gross private domestic investment (GPDI) 
\itemitem{    } \hskip 1.5truein GPDI fixed investment
\itemitem{    } \hskip 1.5truein GPDI change in private inventories
\itemitem{    } \hskip 1.truein Net exports of goods and services 
\itemitem{    } \hskip 1.truein Government consumption expenditures 
                           and gross investment 
\medskip
\itemitem{    } \hskip .5truein www.bea.gov, old NIPA Table 1.9 (in bil.~\$)
\itemitem{    } \hskip 1.truein Gross national product (GNP)
\medskip
\itemitem{    } \hskip .5truein www.bea.gov, old NIPA Table 3.9 (in bil.~\$)
\itemitem{    } \hskip 1.truein Government consumption expenditures 
\medskip
\itemitem{    } \hskip .5truein www.bea.gov, old NIPA Table 5.1 (in bil.~\$)
\itemitem{    } \hskip 1.truein Gross government investment
\medskip
\itemitem{    } \hskip .5truein www.bea.gov, old NIPA Table 3.7 (in bil.~\$)
\itemitem{    } \hskip 1.truein National defense gross investment in structures
\itemitem{    } \hskip 1.truein National defense gross investment in 
                  equipment and software
\medskip 
\itemitem{    } \hskip .5truein www.bea.gov, old NIPA Table 3.5 (in bil.~\$)
\itemitem{    } \hskip 1.truein Federal excise taxes
\itemitem{    } \hskip 1.truein State and local sales taxes
\itemitem{    } \hskip 1.truein State and local other taxes
\medskip 
\itemitem{    } \hskip .5truein www.bea.gov, old NIPA Table 6.2 (in mil.~\$)
\itemitem{    } \hskip 1.truein Compensation of military employees
\medskip 
\itemitem{    } \hskip .5truein www.bea.gov, old NIPA Table 6.3 (in mil.~\$)
\itemitem{    } \hskip 1.truein Wage and salary accruals for military
\bigskip 
\itemitem{\cir} Fixed assets and investments
\itemitem{    } \hskip .5truein www.bea.gov, FA Table 1.1 (in mil.~\$ year-end)
\itemitem{    } \hskip 1.truein Current-cost net stock of fixed assets, 
                 private
\itemitem{    } \hskip 1.truein Current-cost net stock of consumer 
                 durable goods
\medskip 
\itemitem{    } \hskip .5truein www.bea.gov, FA Table 7.1 (in mil.~\$ year-end)
\itemitem{    } \hskip 1.truein Current-cost net stock of fixed assets, 
                 government
\itemitem{    } \hskip 1.truein Current-cost net stock of fixed assets, 
                 government national defense, 
\itemitem{    } \hskip 1.5truein equipment and software
\itemitem{    } \hskip 1.truein Current-cost net stock of fixed assets, 
                 government national defense, 
\itemitem{    } \hskip 1.5truein military facilities
\medskip 
\itemitem{    } \hskip .5truein www.bea.gov, FA Table 1.5 (in mil.~\$)
\itemitem{    } \hskip 1.truein Historical-cost investment in fixed 
                 assets, private
\itemitem{    } \hskip 1.truein Historical-cost investment in consumer 
                 durable goods
\medskip 
\itemitem{    } \hskip .5truein www.bea.gov, FA Table 7.5 (in mil.~\$)
\itemitem{    } \hskip 1.truein Historical-cost investment in fixed assets, 
                 government
\itemitem{    } \hskip 1.truein Historical-cost investment in fixed assets, 
                 government national 
\itemitem{    } \hskip 1.5truein defense, equipment and software
\itemitem{    } \hskip 1.truein Historical-cost investment in fixed assets, 
                 government national 
\itemitem{    } \hskip 1.5truein defense, military facilities
\bigskip 
\itemitem{\cir} Civilian manhours, pre-1947
\itemitem{    } \hskip .5truein Kendrick (1961), Table A-X (millions)
\bigskip 
\itemitem{\cir} Civilian manhours, post-1947
\itemitem{    } \hskip .5truein www.bea.gov, old NIPA Table 6.5 (thousands)
\itemitem{    } \hskip 1.truein Full-time equivalent employees
\itemitem{    } \hskip 1.truein Full-time equivalent employees, military
\medskip 
\itemitem{    } \hskip .5truein www.bea.gov, old NIPA Table 6.8 (thousands)
\itemitem{    } \hskip 1.truein Persons engaged in production 
\itemitem{    } \hskip 1.truein Persons engaged in production, military
\medskip 
\itemitem{    } \hskip .5truein www.bea.gov, old NIPA Table 6.9 (mil.~of hours)
\itemitem{    } \hskip 1.truein Hours worked by full-time and part-time employees

\bigskip
\itemitem{\cir} Population over 16
\itemitem{    } \hskip .5truein {\it Historical Statistics}, Series A6-8
\itemitem{    } \hskip .5truein {\it Economic Report of the President} (2001), 
Table B-34.




\subsection{Appendix Figures}

Figures A1 and A2 are the efficiency and tax wedges, respectively,
over the entire twentieth century.
% In usdata.m: 
% psi=2.24;theta=.35;
% lam=1.016.^(t-1929);
% A=(Ypc./lam)./( (Kpc./lam).^theta .* L.^(1-theta)); 
% taun = 1-(psi/(1-theta)) *(Cpc./Ypc) .*(L./(1-L));


\section{U.S.~Postwar Quarterly Data Measures and Sources}

For the quarterly series, we have total hours data only
after 1959:1.  Since there are no world wars in our postwar
sample, we do not adjust for military compensation.
The main output series that we use for
the postwar series is gross domestic product.  

\subsection{Measures}

\item{\bul} Per-capita output ($y$)
\itemitem{    } \ \ Real GDP 
\itemitem{    } \ \ $-$ Sales tax deflated by the personal consumption 
expenditures (PCE) deflator
\itemitem{    } \ \ $+$ Services from consumer durables (with return = 4\%)
                        deflated by the PCE durable deflator
\itemitem{    } \ \ + Depreciation from consumer durables deflated
                       by the PCE durable deflator
\itemitem{    } \ \ All divided by non-institutional population 16--64
\bigskip
\item{\bul} Per-capita investment ($x$)
\itemitem{    } \ \ Real gross private domestic investment (fixed plus inventories)
\itemitem{    } \ \ + Real government gross investment
\itemitem{    } \ \ + Real personal consumption expenditures on durables
\itemitem{    } \ \ $-$ Sales tax deflated by PCE deflator
                    $\times$ share of durables in PCE 
\itemitem{    } \ \ All divided by non-institutional population 16--64
\bigskip
\item{\bul} Per-capita government ($g$)
\itemitem{    } \ \ Real government consumption
\itemitem{    } \ \ + Real net exports of goods and services
\itemitem{    } \ \ All divided by non-institutional population 16--64
\bigskip
\item{\bul} Per-capita labor input ($l$)
\itemitem{    } \ \ Total hours from the current population survey
\itemitem{    } \ \ + Military hours 
\itemitem{    } \ \ All divided by non-institutional population 16--64

\subsection{Sources}

The specific sources of the data listed above are as follows:
\medskip
\itemitem{\cir} National accounts, post-1947, quarterly
\itemitem{    } \hskip .5truein www.bea.gov, new NIPA Table 1.1.6 
(in bil.~chained 2000 \$)
\itemitem{    } \hskip 1.truein Gross domestic product (GDP)
\itemitem{    } \hskip 1.truein Personal consumption expenditures (PCE)
\itemitem{    } \hskip 1.truein Gross private domestic investment (GPDI) 
\itemitem{    } \hskip 1.truein Net exports of goods and services 
\itemitem{    } \hskip 1.truein Government consumption expenditures 
                           and gross investment 
\medskip
\itemitem{    } \hskip .5truein www.bea.gov, new NIPA Table 1.1.5 (in bil.~\$)
\itemitem{    } \hskip 1.truein PCE durable goods
\itemitem{    } \hskip 1.truein PCE nondurable goods
\itemitem{    } \hskip 1.truein PCE services
\medskip
\itemitem{    } \hskip .5truein www.bea.gov, new NIPA Table 1.1.9 
                        (in 2000 = 100)
\itemitem{    } \hskip 1.truein Deflator, PCE durable goods
\itemitem{    } \hskip 1.truein Deflator, PCE nondurable goods
\itemitem{    } \hskip 1.truein Deflator, PCE services
\medskip
\itemitem{    } \hskip .5truein www.bea.gov, new NIPA Table 3.9.5 (in bil.~\$)
\itemitem{    } \hskip 1.truein Government consumption expenditures 
\medskip 
\itemitem{    } \hskip .5truein www.bea.gov, new NIPA Table 3.2 (in bil.~\$)
\itemitem{    } \hskip 1.truein Federal excise taxes
\medskip 
\itemitem{    } \hskip .5truein www.bea.gov, new NIPA Table 3.3 (in bil.~\$)
\itemitem{    } \hskip 1.truein State and local sales taxes
\itemitem{    } \hskip 1.truein State and local other taxes
\bigskip 
\itemitem{\cir} Flow of funds accounts, post-1952, quarterly
\itemitem{    } \hskip .5truein www.federalreserve.gov, Flow Table 10 (in mil.~\$)
\itemitem{    } \hskip 1.truein Consumption of fixed capital, consumer durables
\medskip 
\itemitem{    } \hskip .5truein www.federalreserve.gov, Level Table 100 (in bil.~\$)
\itemitem{    } \hskip 1.truein Current-cost net stock of consumer durables
\medskip 
\itemitem{\cir} Hours and population, post-1959, quarterly
\itemitem{    } \hskip .5truein Prescott, Ueberfeldt, and Cociuba (2005), Hours.xls
\itemitem{    } \hskip 1.truein Non-institutional hours from the current
population survey
\itemitem{    } \hskip 1.truein Non-institutional population, ages 16--64
\bigskip




\Appendix{C}{Additional Material Not Reported in Text}


\noindent
In this Appendix, we describe further details on 
results that were mentioned only briefly in the main
text of the paper.

\section{Results for Government Consumption Wedge}

Figures A3 and A4 show results with varying
government consumption wedges in the benchmark model
for the Great Depression and the 1982 recession, respectively.
The results for extensions of the benchmark
model---with variable capital utilization and 
with adjustment costs---are almost identical
to Figure A3 and therefore not shown.

\section{Model with Maximum Investment Wedge}

In the paper, we reported on results for the {\it Model with Maximum
Investment Wedge}.  In this case, the investment
wedge was chosen to be as large as it needed to 
be in order for investment in the model and data
to line up.  In Figure A5, we display the investment
series---both model and data---along with consumption
and output.  (Government consumption in the model
is set equal to the 1929 level, but there is little
change in the data over the period 1929--1939.)
Notice that consumption in the model rises significantly
as investment falls implying a consumption anomaly.

\section{Wedges for the Capital Utilization Model}

Figure A6 shows all wedges for the model with 
variable capital utilization.  In Figure 9 of the 
paper we displayed only the efficiency wedge.
Here, we provide the exact analogue of Figure 1
with all wedges but the government consumption wedge.

\section{Results for an Alternative Investment Wedge}

Figure A7 shows that our accounting procedure is
not sensitive to choosing $\tau_k$ as the investment
wedge rather than $\tau_x$.  This figure shows the 
model's prediction for output based on the two prototype
models described in Sections A.1 and A.2.4.  
For both, we set the adjustment costs to 
be extreme ($\eta=1$) so that we could compare
the results to Christiano and Davis (2006).  

The line marked {\it Prediction With $\tau_x$ Wedge and Extreme
Costs}
in Figure A7
is the same as that in Figure 14 in our paper.
The line marked {\it Prediction With $\tau_k$ Wedge
and Extreme Costs}
is what Christiano and Davis (2006)
would find if they were to apply the same accounting
procedure that we do.
Here, we use 
Christiano and Davis' (2006) 
estimates for the 
stochastic process, namely,
$$\EQNalign{
  s_{t+1} = \left[\matrix{-.0100\cr
                           .0279\cr
                           .0010\cr
                           -.0002}\right]
          & + \left[\matrix{.9026 & .0296  & -0.0009 & -.0007\cr
                            .0766 & .9350  & .4492   & .0078\cr
                            .0171 & -.0013 & .9333   & .0000\cr
                            .0079 & .0050  & .0056   & 0.9976\cr
                    }\right] s_{t}\cr
\noalign{\bigskip}
          &
          + \left[\matrix{-.0130 & 0 & 0 & 0\cr
                          -.0013 & .0065 &  0  & 0\cr
                           .0040 & .0007 & -.0017 &  0\cr
                           .0017 & .0077 & -.0069 & .0137\cr
                         }\right] \epsilon_{s,t+1}.
}
$$

In Figure A8, we redisplay the line marked {\it Prediction With
$\tau_k$ Wedge and Extreme Costs} from Figure A7.  In Figure A8, 
we label it
{\it Prediction Using a Theoretically-Consistent Methodology}.
We want to compare this result with the result that
Christiano and Davis (2006) actually report, what we
refer to as {\it Prediction Using an Alternative Methodology.}

As we noted earlier,
our propositions distinguish between the direct
effect and the forecasting effect of fluctuations
in wedges.   This turned out not to be 
quantitatively important for the 
prototype model with a $\tau_x$ investment
wedge.\note{%
Even so, we redid all of our numerical results 
applying the accounting procedure that is consistent with 
our propositions.}
However, it is for the $\tau_k$ prototype model
as Figure A8 shows.
In fact, because Christiano and Davis (2006)
are not consistently applying
the propositions, they find that the investment
wedge accounts for one-half of the downfall in output.
In fact, this wedge accounts for only about 
one-fifth of the downfall
with the caveat that, even then, adjustment costs have to be
extreme.

In summary, our main conclusion---that the investment wedge
plays a decidedly tertiary role---is not sensitive to
using the alternative prototype model proposed by Christiano
and Davis (2006).

\section{Proof of Propositions}

Here, we provide a proof of Proposition 1.  The
logic behind the other proofs is the same.

The first-order conditions for the gross output
problem are 
$p_{1t}q_{1t}=\gamma q_t$ 
and
$p_{2t}q_{2t}=(1-\gamma) q_t$.
The first-order conditions for the sector $i$ gross
output problem are
$\theta p_{it}q_{it}/m_{it}=1+\tau_{it}$
and $(1-\theta)p_{it}q_{it}/z_{it}=v_t$.
Thus, 
$m_{1t}=\theta\gamma q_t/(1+\tau_{1t})$
and 
$m_{2t}=\theta\gamma q_t/(1+\tau_{2t})$
so 
$m_t=\theta q_t[\gamma/(1+\tau_{1t})+(1-\gamma)/(1+\tau_{2t})]$.
The first-order conditions also imply $z_{1t}=\gamma z_t$,
$z_{2t}=(1-\gamma)z_t$, and $v_t=(1-\theta)q_t/z_t$.
Hence,
$$
 y_t = q_t - m_t = [1-\theta(a_{1t}+a_{2t})]q_t,
\EQN A1
$$
and substituting the expressions for $m_{it}$ and $z_{it}$
into $q_t=(m_{1t}^\theta z_{1t}^{1-\theta})^\gamma
          (m_{2t}^\theta z_{2t}^{1-\theta })^{1-\gamma}$ 
and manipulating gives
$$
q_{t}=\kappa (a_{1t}^{1-\gamma }a_{2t}^{\gamma })^{\theta \over 1-\theta} z_t.  
\EQN {A2}
$$
Combining the expressions for $y_t$ and $q_t$ and using 
$z_t=F(k_t,l_t)$ gives the expression for $A_t$.

To verify our expression for $(1-\tau _{lt})$, compare the first-order
condition $v_{t}F_{lt}=w_{t}$ from the composite goods producer in the
detailed economy to the first-order condition 
$(1-\tau_{lt})A_tF_{lt}=w_t$
in the prototype economy to note that 
$(1-\tau_{lt})=v_t/A_t$
and use $v_t=(1-\theta )q_t/z$, \Ep{A2}, 
and the expression for $A_t$.
The derivation of the expression for $(1-\tau _{kt})$ is analogous. 

\section{MLE Estimates for Alternative Models}

In Table A1, we report the parameter estimates for
the model with variable capital utilization  and annual data
over the period 1901--1940.
In Tables A2 and A3, we report the estimates for
the model with adjustment costs, at the BGG level
and four times the BGG level, based on annual data
for 1901--1940.  The quarterly estimates for the 
adjustment-cost models are reported in Tables A4 and
A5.  These estimates are based on quarterly data
for the period 1959:1--2004:3.   

The hillclimbing procedure that we used had no problems
finding maxima (at least locally) with the annual
dataset during the 1901--1940 period.  Estimation in 
the postwar period was more difficult, and the hillclimbing
procedure oftentimes had difficulty finding higher points
on the likelihood surface.  Thus, in all cases for the postwar,
we initialized guesses using estimates based on 
annual data for  the period 1901--2000, which were
converted to quarterly estimates.
We also perturbed the parameter estimates (after
the hillclimbing routine could not find a higher point)
many times in search of higher points.  
We did not find our results to be sensitive to
the estimates of the stochastic process.


\newpage
\centerline{\twelvepoint\bf  References}
\bigskip
\bigskip
{
                                                                                
\baselineskip=13pt
\parskip = 5pt
\parindent= -.25truein
\leftskip = .25truein
                                                                                
\vskip -10pt

Christiano, Lawrence and Joshua Davis. 2006.
Two Flaws in Business Cycle Accounting.
Manuscript, Northwestern University.

Federal Reserve Board of Governors.
{\it Flow of Funds Accounts}. 1945--2005.
Washington, DC: Federal Reserve Board.
(Contains flow and level tables, the latest of which
are available at www.federalreserve.gov.)

Kendrick, John W. 1961. 
{\it Productivity Trends in the United States}.
Princeton, NJ: Princeton University Press.

Prescott, Edward C., Alexander Ueberfeldt, and Simona Cociuba. 2005.
U.S.~Hours and Productivity Behavior Using CPS Hours
Worked Data: 1959-I to 2004-IV.
Manuscript, Federal Reserve Bank of Minneapolis.

Tauchen, George. 1986.
Finite State Markov-Chain Approximations to Univariate and 
 Vector Autoregressions.
{\it Economic Letters}, 20(2): 177--181.   

U.S.~Department of Commerce. 
Bureau of the Census. 1975.
{\it Historical Statistics of the United States, Colonial to 1970}.
Washington, DC: U.S. Government Printing Office.

U.S. Department of Commerce. 
Bureau of Economic Analysis. 1929--2005.
{\it Survey of Current Business}.
Washington, DC: U.S. Government Printing Office.
(Contains NIPA tables and fixed asset tables, the latest of which
               are available at www.bea.gov.)

U.S.~Executive Office. Council of Economic Advisers. 2001.
{\it Economic Report of the President.}
Washington, DC: U.S. Government Printing Office.

}

\newpage
{\ }
\vskip -1truein
\def\epsfsize#1#2{.6#1}
\centerline{\epsffile{/home/erm/papers/econ2/results/figures/FigureA1.ps}}
                                                                                
\newpage
{\ }
\vskip -1truein
\def\epsfsize#1#2{.6#1}
\centerline{\epsffile{/home/erm/papers/econ2/results/figures/FigureA2.ps}}

\newpage
{\ }
\vskip -1truein
\def\epsfsize#1#2{.6#1}
\centerline{\epsffile{/home/erm/papers/econ2/results/figures/FigureA3.ps}}

\newpage
{\ }
\vskip -1truein
\def\epsfsize#1#2{.6#1}
\centerline{\epsffile{/home/erm/papers/econ2/results/figures/FigureA4.ps}}

\newpage
{\ }
\vskip -1truein
\def\epsfsize#1#2{.6#1}
\centerline{\epsffile{/home/erm/papers/econ2/results/figures/FigureA5.ps}}

\newpage
{\ }
\vskip -1truein
\def\epsfsize#1#2{.6#1}
\centerline{\epsffile{/home/erm/papers/econ2/results/figures/FigureA6.ps}}

\newpage
{\ }
\vskip -1truein
\def\epsfsize#1#2{.6#1}
\centerline{\epsffile{/home/erm/papers/econ2/results/figures/FigureA7.ps}}

\newpage
{\ }
\vskip -1truein
\def\epsfsize#1#2{.6#1}
\centerline{\epsffile{/home/erm/papers/econ2/results/figures/FigureA8.ps}}

\newpage
\centerline{\bf Table A1}
\centerline{\bf Parameters 
for the Model with Variable Capital Utilization}
\centerline{\bf  Annual Data, 1901--1940 }
\medskip
\hrule
\vskip 1pt
\hrule
$$\eqalign{
%\noalign{\smallskip}
%&\hskip -.1truein {\bf Annual\ Data,\ 1901\!-\!1940}\cr
\noalign{\bigskip}
  &\hskip .75truein P =\ \left[\matrix{
              {.666}\atop {(\scriptscriptstyle{.415,.820})} 
                   & {.171}\atop{(\scriptscriptstyle{.0626,.278})} 
                   & {-.192}\atop{(\scriptscriptstyle{-.390,.0205})} 
                   & 0 \cr
\noalign{\vskip 3pt}
              {-.178}\atop{(\scriptscriptstyle{-.342,.0122})} 
                   & {1.08}\atop{(\scriptscriptstyle{.926,1.10})} 
                   & {.285}\atop{(\scriptscriptstyle{.105,.500})} 
                   & 0 \cr
\noalign{\vskip 3pt}
             {-.0402}\atop{(\scriptscriptstyle{-.434,.367})} 
                   & {.0391}\atop{(\scriptscriptstyle{-.208,.317})} 
                   & {.108}\atop{(\scriptscriptstyle{-.267,.289})} 
                   & 0\cr
\noalign{\vskip 3pt}
                   0
                   & 0 
                   & 0
                   & {.744}\atop{(\scriptscriptstyle{.468,.782})} \cr
        }
        \right]\cr
\noalign{\bigskip}
\noalign{\bigskip}
  &\hskip .75truein Q =\ \left[\matrix{
            {.0325}\atop{(\scriptscriptstyle{.0249,.0373})} 
                & 0 
                & 0 
                & 0 \cr
\noalign{\vskip 3pt}
            {.0109}\atop{(\scriptscriptstyle{.00414,.0207})} 
                & {.0342}\atop{(\scriptscriptstyle{.0247,.0389})} 
                & 0 
                & 0 \cr
\noalign{\vskip 3pt}
            {.00971}\atop{(\scriptscriptstyle{-.00976,.0265})} 
                & {-.000895}\atop{(\scriptscriptstyle{-.0153,.0160})} 
                & {.0375}\atop{(\scriptscriptstyle{.0187,.0501})} 
                & 0 \cr
\noalign{\vskip 3pt}
                0
                & 0
                & 0
                & {.222}\atop{(\scriptscriptstyle{.147,.281})} \cr
           }
        \right] \cr
\noalign{\bigskip}
& {\rm Mean}(s)
  = [.744\,(.708,.787),\ .229\,(.175,.288),\ 
                      .282\,(.227,.333),\ -2.78\,(-2.94,-2.53)]\cr
}
$$
\hrule
\vskip 1pt
\hrule
\medskip
{ \baselineskip 10pt

\item{} {\eightpoint Note: Parameters were estimated using maximum
likelihood with data on output, labor, investment, and government
consumption.
To ensure stationarity, we added a penalty term to the likelihood
function proportional to $\max(|\lambda_{\rm max}|-.995,0)^2$,
where $\lambda_{\rm max}$ is the maximal eigenvalue of $P$.
Numbers in parentheses are 90\% confidence
intervals for a bootstrapped distribution with 500 replications.}
                                                                                
}
                                                                                


\newpage
\centerline{\bf Table A2}
\centerline{\bf Parameters 
for the Model with Adjustment Costs at the BGG Level}
\centerline{\bf  Annual Data, 1901--1940 }
\medskip
\hrule
\vskip 1pt
\hrule
$$\eqalign{
%\noalign{\smallskip}
%&\hskip -.1truein {\bf Annual\ Data,\ 1901\!-\!1940}\cr
\noalign{\bigskip}
  &\hskip .75truein P =\ \left[\matrix{
              {.564}\atop {(\scriptscriptstyle{.225,.767})} 
                   & {.0898}\atop{(\scriptscriptstyle{.0213,.192})} 
                   & {-.190}\atop{(\scriptscriptstyle{-.405,.0241})} 
                   & 0 \cr
\noalign{\vskip 3pt}
              {.0470}\atop{(\scriptscriptstyle{-.246,.317})} 
                   & {.995}\atop{(\scriptscriptstyle{.879,1.07})} 
                   & {.219}\atop{(\scriptscriptstyle{.00612,.455})} 
                   & 0 \cr
\noalign{\vskip 3pt}
             {-.378}\atop{(\scriptscriptstyle{-.919,.00349})} 
                   & {.0545}\atop{(\scriptscriptstyle{-.190,.225})} 
                   & {.337}\atop{(\scriptscriptstyle{-.0208,.601})} 
                   & 0\cr
\noalign{\vskip 3pt}
                   0
                   & 0 
                   & 0
                   & {.766}\atop{(\scriptscriptstyle{.401,.842})} \cr
        }
        \right]\cr
\noalign{\bigskip}
\noalign{\bigskip}
  &\hskip .75truein Q =\ \left[\matrix{
            {.0567}\atop{(\scriptscriptstyle{.0431,.0641})} 
                & 0 
                & 0 
                & 0 \cr
\noalign{\vskip 3pt}
            {-.00425}\atop{(\scriptscriptstyle{-.0210,.0114})} 
                & {.0554}\atop{(\scriptscriptstyle{.0380,.0638})} 
                & 0 
                & 0 \cr
\noalign{\vskip 3pt}
            {-.0432}\atop{(\scriptscriptstyle{-.0839,-.0103})} 
                & {.0206}\atop{(\scriptscriptstyle{-.00282,.0409})} 
                & {.0768}\atop{(\scriptscriptstyle{.0453,.0922})} 
                & 0 \cr
\noalign{\vskip 3pt}
                0
                & 0
                & 0
                & {.221}\atop{(\scriptscriptstyle{.143,.271})} \cr
           }
        \right] \cr
\noalign{\bigskip}
& {\rm Mean}(s)
  = [.537\,(.499,.582),\ -.196\,(-.277,-.0941),\ 
                      .291\,(.191,.393),\ -2.80\,(-2.97,-2.53)]\cr
}
$$
\hrule
\vskip 1pt
\hrule
\medskip
{ \baselineskip 10pt
                                                                                
\item{} {\eightpoint See footnotes to Table A1.}

}

%\bigskip


\newpage
\centerline{\bf Table A3}
\centerline{\bf Parameters 
for the Model with Adjustment Costs at 4 times the BGG Level}
\centerline{\bf  Annual Data, 1901--1940 }
\medskip
\hrule
\vskip 1pt
\hrule
$$\eqalign{
%\noalign{\smallskip}
%&\hskip -.1truein {\bf Annual\ Data,\ 1901\!-\!1940}\cr
\noalign{\bigskip}
  &\hskip .75truein P =\ \left[\matrix{
              {.432}\atop {(\scriptscriptstyle{-.0446,.625})} 
                   & {.121}\atop{(\scriptscriptstyle{.0476,.276})} 
                   & {-.0866}\atop{(\scriptscriptstyle{-.191,-.0167})} 
                   & 0 \cr
\noalign{\vskip 3pt}
              {.178}\atop{(\scriptscriptstyle{-.160,.603})} 
                   & {.963}\atop{(\scriptscriptstyle{.796,1.03})} 
                   & {.0938}\atop{(\scriptscriptstyle{.0169,.196})} 
                   & 0 \cr
\noalign{\vskip 3pt}
             {-.659}\atop{(\scriptscriptstyle{-1.00,1.00})} 
                   & {.112}\atop{(\scriptscriptstyle{-.419,.331})} 
                   & {.560}\atop{(\scriptscriptstyle{.264,.942})} 
                   & 0\cr
\noalign{\vskip 3pt}
                   0
                   & 0 
                   & 0
                   & {.783}\atop{(\scriptscriptstyle{.433,.874})} \cr
        }
        \right]\cr
\noalign{\bigskip}
\noalign{\bigskip}
  &\hskip .75truein Q =\ \left[\matrix{
            {.0561}\atop{(\scriptscriptstyle{.0422,.0635})} 
                & 0 
                & 0 
                & 0 \cr
\noalign{\vskip 3pt}
            {-.00327}\atop{(\scriptscriptstyle{-.0209,.0143})} 
                & {.0553}\atop{(\scriptscriptstyle{.0381,.0637})} 
                & 0 
                & 0 \cr
\noalign{\vskip 3pt}
            {-.176}\atop{(\scriptscriptstyle{-.252,-.0896})} 
                & {.0811}\atop{(\scriptscriptstyle{.0247,.137})} 
                & {.189}\atop{(\scriptscriptstyle{.119,.216})} 
                & 0 \cr
\noalign{\vskip 3pt}
                0
                & 0
                & 0
                & {.221}\atop{(\scriptscriptstyle{.143,.271})} \cr
           }
        \right] \cr
\noalign{\bigskip}
& {\rm Mean}(s)
  = [.536\,(.497,.577),\ -.200\,(-.274,-.0962),\ 
                      .292\,(.163,.432),\ -2.81\,(-3.00,-2.54)]\cr
}
$$
\hrule
\vskip 1pt
\hrule
\medskip
{ \baselineskip 10pt
                                                                                
\item{} {\eightpoint See footnotes to Table A1.}

}

% A4:
%  -0.02593003143300  -0.02606184437574  -0.02582513727466
%   0.32528078240235   0.32509529754454   0.32542637467004
%   0.47640296988723   0.47615437754812   0.47644392281429
%  -1.52572392233495  -1.52585919158196  -1.52539951477464

%   0.98814714427224   0.98778603498219   0.98845592338782
%  -0.00393620952383  -0.00427457976845  -0.00381089232356
%   0.00739476978480   0.00717194283672   0.00744481871123
%  -0.00154763964840  -0.00189988543225  -0.00113335210194
%   0.02347540555172   0.02317549494312   0.02389753009218
%   0.99522631431640   0.99340031776425   0.99525568484684
%  -0.07813326483521  -0.07822096658493  -0.07795594109085
%  -0.01292726594313  -0.01337987268383  -0.01282330502371
%   0.05213951101474   0.05114164861296   0.05213995058101
%   0.03785132235434   0.03742531716084   0.03809635811273
%   0.87933464895897   0.87918748861040   0.87989414805935
%   0.04800537603439   0.04744605247261   0.04817189345267
%  -0.05623786887083  -0.05705311211362  -0.05623786812953
%  -0.02974898581502  -0.03031793563749  -0.02967043228482
%   0.14558343713444   0.14543892403670   0.14650252963377
%   0.96500911999557   0.96418535249112   0.96507887177841
%   0.01143591934237   0.01086446554766   0.01199525827391
%   0.00132263476266   0.00072272505497   0.00198247543326
%  -0.00924790891037  -0.00987114817555  -0.00863183594142
%  -0.00038806619786  -0.00091033505320   0.00025218732691
%   0.00640541868455   0.00575960701346   0.00701166479057
%   0.00264329092744   0.00175280222628   0.00324958905672
%   0.00619191723283   0.00540291740960   0.00681142882766
%   0.02206456706386   0.02127018943454   0.02266601525908
%   0.01348046236832   0.01262397697590   0.01413874049684
%   0.00697830381845   0.00626451974025   0.00761658601290


\newpage
\centerline{\bf Table A4}
\centerline{\bf Parameters 
for the Model with Adjustment Costs at the BGG Level}
\centerline{\bf  Quarterly Data, 1959:1--2004:3 }
\medskip
\hrule
\vskip 1pt
\hrule
$$\eqalign{
%\noalign{\smallskip}
%&\hskip -.1truein {\bf Quarterly\ Data,\ 1959\!:1\!-\!2004\!:3}\cr
\noalign{\bigskip}
  &\hskip .75truein P =\ \left[\matrix{
              {.988}\atop {(\scriptscriptstyle{.988,.988})} 
                   & {.0235}\atop{(\scriptscriptstyle{.0232,.0239})} 
                   & {.0521}\atop{(\scriptscriptstyle{.0511,.0521})} 
                   & {-.0562}\atop{(\scriptscriptstyle{-.0571,-.0562})} \cr
\noalign{\vskip 3pt}
              {-.00394}\atop{(\scriptscriptstyle{-.00427,-.00381})} 
                   & {.995}\atop{(\scriptscriptstyle{.926,1.10})} 
                   & {.0379}\atop{(\scriptscriptstyle{.0374,.0381})} 
                   & {-.0297}\atop{(\scriptscriptstyle{-.0303,-.0297})} \cr
\noalign{\vskip 3pt}
             {.00739}\atop{(\scriptscriptstyle{.00717,.00744})} 
                   & {-.0781}\atop{(\scriptscriptstyle{-.0782,-.0780})} 
                   & {.879}\atop{(\scriptscriptstyle{.879,.880})} 
                   & {.146}\atop{(\scriptscriptstyle{.145,.147})} \cr
\noalign{\vskip 3pt}
             {-.00155}\atop{(\scriptscriptstyle{-.00190,-.00113})} 
                   & {-.0129}\atop{(\scriptscriptstyle{-.0134,-.0128})} 
                   & {.0480}\atop{(\scriptscriptstyle{.0474,.0482})} 
                   & {.965}\atop{(\scriptscriptstyle{.964,.965})} \cr
        }
        \right]\cr
\noalign{\bigskip}
\noalign{\bigskip}
  &\hskip .75truein Q =\ \left[\matrix{
            {.0114}\atop{(\scriptscriptstyle{.0109,.0120})} 
                & 0 
                & 0 
                & 0 \cr
\noalign{\vskip 3pt}
            {.00132}\atop{(\scriptscriptstyle{.000723,.00198})} 
                & {.00641}\atop{(\scriptscriptstyle{.00576,.00701})} 
                & 0 
                & 0 \cr
\noalign{\vskip 3pt}
            {-.00925}\atop{(\scriptscriptstyle{-.00987,-.00863})} 
                & {.00264}\atop{(\scriptscriptstyle{.00175,.00325})} 
                & {.0221}\atop{(\scriptscriptstyle{.0213,.0227})} 
                & 0 \cr
\noalign{\vskip 3pt}
            {-.000388}\atop{(\scriptscriptstyle{-.000910,.000252})} 
                & {.00619}\atop{(\scriptscriptstyle{.00540,.00681})} 
                & {.0135}\atop{(\scriptscriptstyle{.0126,.0141})} 
                & {.00698}\atop{(\scriptscriptstyle{.00626,.00762})} \cr
           }
        \right] \cr
\noalign{\bigskip}
& {\rm Mean}(s)
  = [-.0239\,(-.0261,-.0258),\ .325\,(.325,.325),\ 
                      .476\,(.476,.476),\ -1.53\,(-1.53,-1.53)]\cr
}
$$
\hrule
\vskip 1pt
\hrule
\medskip
{ \baselineskip 10pt
                                                                                
\item{} {\eightpoint See footnotes to Table A1.}

}

% A5:
%  -0.03868373032989  -0.03868401127952  -0.03868352091685
%   0.32533115894097   0.32533076342042   0.32533131262676
%   0.42386985342091   0.42386965571428   0.42386998980220
%  -1.53913101241965  -1.53913133840198  -1.53913081793379
%   1.01420535776153   1.01420536134203   1.01421293421703
%   0.00700407325041   0.00699903050713   0.00700407160174
%  -0.10552749200191  -0.10552749181397  -0.10552596070579
%  -0.03794064209269  -0.03794210579936  -0.03794064228212
%  -0.04443387632167  -0.04444041116046  -0.04443387402647
%   0.95319275816302   0.95317713255821   0.95319275371482
%   0.10009893037604   0.10009893116483   0.10010158752171
%  -0.01576973857835  -0.01577213732362  -0.01576973861925
%   0.03612391751414   0.03609963495900   0.03612391601289
%   0.01732329566843   0.01728774729174   0.01732328745167
%   0.82961577117656   0.82961577371188   0.82962010969731
%  -0.03018624159039  -0.03019078375273  -0.03018623642559
%  -0.01684629253098  -0.01686990218482  -0.01684629390460
%   0.00403375819985   0.00399873550270   0.00403375009292
%   0.13978989734296   0.13978989983818   0.13979429147891
%   1.04444661366546   1.04444145665586   1.04444662698571
%   0.01137088090910   0.01136788622594   0.01137487629863
%   0.00148119047057   0.00147671763890   0.00148560446021
%  -0.02717953211122  -0.02718032511022  -0.02717852164585
%  -0.00002984273316  -0.00003167746727  -0.00002871809365
%   0.00626157343799   0.00625767287506   0.00626549162188
%   0.01220319595520   0.01220215496172   0.01220475121525
%   0.00589079806760   0.00588843899562   0.00589209546809
%   0.02431630028343   0.02431467690438   0.02431822285346
%   0.00929368330591   0.00929162679219   0.00929509621744
%   0.01154144364817   0.01153961477381   0.01154401874170

\newpage
\centerline{\bf Table A5}
\centerline{\bf Parameters 
for the Model with Adjustment Costs at 4 times the BGG Level}
\centerline{\bf  Quarterly Data, 1959:1--2004:3 }
\medskip
\hrule
\vskip 1pt
\hrule
$$\eqalign{
%\noalign{\smallskip}
%&\hskip -.1truein {\bf Quarterly\ Data,\ 1959\!:1\!-\!2004\!:3}\cr
\noalign{\bigskip}
  &\hskip .75truein P =\ \left[\matrix{
              {1.014}\atop {(\scriptscriptstyle{1.014,1.014})} 
                   & {-.0444}\atop{(\scriptscriptstyle{-.0444,-.0444})} 
                   & {.0361}\atop{(\scriptscriptstyle{.0361,.0361})} 
                   & {-.0168}\atop{(\scriptscriptstyle{-.0168,-.0168})} \cr
\noalign{\vskip 3pt}
              {.00700}\atop{(\scriptscriptstyle{.00700,.00700})} 
                   & {.953}\atop{(\scriptscriptstyle{.953,.953})} 
                   & {.0173}\atop{(\scriptscriptstyle{.0173,.0173})} 
                   & {.00403}\atop{(\scriptscriptstyle{.00400,.00403})} \cr
\noalign{\vskip 3pt}
             {-.106}\atop{(\scriptscriptstyle{-.106,.-.106})} 
                   & {.100}\atop{(\scriptscriptstyle{.100,.100})} 
                   & {.830}\atop{(\scriptscriptstyle{.830,.830})} 
                   & {.140}\atop{(\scriptscriptstyle{.140,.140})} \cr
\noalign{\vskip 3pt}
             {-.0379}\atop{(\scriptscriptstyle{-.0379,-.0379})} 
                   & {-.0158}\atop{(\scriptscriptstyle{-.0158,-.0158})} 
                   & {-.0302}\atop{(\scriptscriptstyle{-.0302,-.0302})} 
                   & {1.044}\atop{(\scriptscriptstyle{1.044,1.044})} \cr
        }
        \right]\cr
\noalign{\bigskip}
\noalign{\bigskip}
  &\hskip .75truein Q =\ \left[\matrix{
            {.0114}\atop{(\scriptscriptstyle{.0114,.0114})} 
                & 0 
                & 0 
                & 0 \cr
\noalign{\vskip 3pt}
            {.00148}\atop{(\scriptscriptstyle{.00148,.00149})} 
                & {.00626}\atop{(\scriptscriptstyle{.00626,.00627})} 
                & 0 
                & 0 \cr
\noalign{\vskip 3pt}
            {-.0272}\atop{(\scriptscriptstyle{-.0272,-.0272})} 
                & {.0122}\atop{(\scriptscriptstyle{.0122,.0122})} 
                & {.0243}\atop{(\scriptscriptstyle{.0243,.0243})} 
                & 0 \cr
\noalign{\vskip 3pt}
            {-.0000298}\atop{(\scriptscriptstyle{-.0000317,-.0000287})} 
                & {.00589}\atop{(\scriptscriptstyle{.00589,.00589})} 
                & {.00929}\atop{(\scriptscriptstyle{.00929,.00929})} 
                & {.0115}\atop{(\scriptscriptstyle{.0115,.0115})} \cr
           }
        \right] \cr
\noalign{\bigskip}
& {\rm Mean}(s)
  = [-.0387\,(-.0387,-.0387),\ .325\,(.325,.325),\ 
                      .424\,(.424,.424),\ -1.54\,(-1.54,-1.54)]\cr
}
$$
\hrule
\vskip 1pt
\hrule
\medskip
{ \baselineskip 10pt
                                                                                
\item{} {\eightpoint See footnotes to Table A1.}

}


\bye

